%%%%%%%%%%%%%%%%%%%%%%%%%%%%%%%%%%%%%%%%%%%%%%%%%%%%%%%%%%%%%%%%%%%%%%%%%%%%%%%%%%%%%
%% Modelo adaptado de "Modelo de TCC do Curso de Engenharia de Controle e Automação da UFSC - Campus Blumenau" em mar/2022.
%%% Adaptado em nov/2025 por Tatiana Pazelli
%%%%%%%%%%%%%%%%%%%%%%%%%%%%%%%%%%%%%%%%%%%%%%%%%%%%%%%%%%%%%%%%%%%%%%%%%%%%%%%%%%%%%

\documentclass[12pt, openany, oneside, a4paper, english, brazil]{abntex2}   % Padrão Abntex2

\input{style/packages.tex} % resto dos pacotes Latex

%\usepackage[english]{style/UFSCar-DEE-Monograph} % personalização da ABNTEX2 para escrita em inglês
\usepackage[portuguese]{style/UFSCar-DEE-Monograph} % personalização da ABNTEX2 para escrita em português


%---------------------------------------------------------------
%--------- DADOS BÁSICOS DO TRABALHO (Preencher todos) ---------
%---------------------------------------------------------------

\titulo{Previsão de Emissões de CO$_2$ de Usinas Termelétricas Nacionais utilizando Redes Neurais Artificiais}   % Título do Trabalho em Português
\tituloingles{CO$_2$ Emissions Forecasting for National Thermoelectric Power Plants using Artificial Neural Networks}   % Título do Trabalho em Inglês

\palavraschave{Emissões de CO$_2$; Redes Neurais Artificiais; Termelétricas; Calibração; Regularização L2.}  % 3-6 palavras-chave
\keywords{CO$_2$ Emissions; Artificial Neural Networks; Thermal Power Plants; Calibration; L2 Regularization.}   % 3-6 palavras-chave para resumo em inglês

\DEELautor{Giovanni Grizante}{Aurelio} % Aluno autor do trabalho

\curso{Engenharia Elétrica}												
\orientador{Prof. Dr. Guilherme Guimarães Lage}  % Professor Orientador membro da banca 
\orientadoruni{Universidade Federal de São Carlos} % Universidade de origem do professor orientador

\membrob{Prof. Dr. XXX1}   % Professor membro da banca
\membrobuni{Universidade XXX1}   % Universidade de origem do professor membro da banca 
\membroc{Prof. Dr. XXX2}   % Professor membro da banca
\membrocuni{Universidade XXX2}   % Universidade de origem do professor membro da banca 

\local{São Carlos}   % Cidade da instituição
\data{\the\year}        % Ano de realização


%--------------------------------------------------------------------------------------------

\begin{document}
% Seleção automática do idioma, não alterar
\ifx\isenglish\undefined
    \selectlanguage{brazil}
\else
    \selectlanguage{english}
\fi

% Texto da Dedicatória (OBS: dedicatória é opcional)
% Caso não queira inserir, deixe em branco ( \dedicatoria{} )
\DEELdedicatoria{Dedico este trabalho a todos que me ajudaram a superar os desafios, em especial, aos meus colegas da turma Elétrica 019. Sem vocês, tudo seria mais difícil.}  
    
% Texto dos Agradecimentos (OBS: Agradecimentos é opcional)
% Caso não queira inserir, deixe em branco ( \agradecimentos{} )
\DEELagradecimentos{Ao meu orientador, Prof. Dr. Guilherme Lage, minha eterna gratidão pela paciência, pelas discussões enriquecedoras e por me guiar com leveza mesmo nos momentos mais difíceis.

Aos meus pais, por acreditarem nos meus estudos desde o primeiro dia. Agradeço pelos sacrifícios silenciosos, pelo colo nos dias de cansaço e por nunca medirem esforços para me ver feliz. 

Ao meu amor, Bruna, por compreender minhas ausências, por ouvir meus ensaios cansativos e por segurar minha mão cada vez que pensei em desistir. Você foi meu porto seguro. 

Aos meus amigos de graduação, especialmente aos do `Grupo da Família', obrigado pelas horas na biblioteca, ensinamentos (e cervejas) compartilhados, e por transformarem um processo solitário em uma jornada conjunta.

À CNPq, pelo apoio financeiro concedido por meio da bolsa de iniciação científica, sem o qual esta pesquisa não seria viável.} 
    
                                  
% Texto da Epígrafe (OBS: Epígrafe é opcional)
% Caso não queira inserir, deixe em branco ( \epigrafe{} )
\DEELepigrafe{\textit{"Nós sempre nos definimos pela capacidade de superar o impossível."} \\ (Joseph Cooper)} 
    

% Texto do Resumo (OBS: Resumo é Obrigatório)
\DEELresumo{Embora a maior parte da energia elétrica no Brasil provenha de usinas hidroelétricas, 
as termoelétricas são fundamentais para a segurança energética nacional. No entanto, com as mudanças 
climáticas causadas pelas emissões de CO$_2$ de diversos setores econômicos, é crucial equilibrar 
essa segurança com as emissões de CO$_2$ no setor elétrico brasileiro. Este trabalho propõe estimar 
as emissões horárias de CO$_2$ equivalente (CO$_2$-e) das termoelétricas brasileiras usando redes 
neurais para apoiar o planejamento energético e ambiental do Sistema Interligado Nacional. 
Enquanto a literatura baseia a modelagem dessas emissões em funções específicas para cada usina, 
as emissões no Brasil são inventariadas pela Eletrobras e pelo Instituto de Energia e Meio Ambiente 
em uma base anual. Para contornar isso, este estudo discretiza as emissões anuais de cada usina em 
estimativas horárias, utilizando regressão não linear em função de dados disponibilizados pelo 
Operador Nacional do Sistema Elétrico. Com isso, é possível generalizar as estimativas de emissões 
horárias das termoelétricas brasileiras. Este trabalho inclui análises para 39 usinas termelétricas 
com diferentes capacidades de geração, ciclos de operação e combustíveis, fornecendo \textit{insights} 
para o planejamento energético, elétrico e ambiental do país.} 
     
% Texto do Abstract (OBS: Abstract é Obrigatório)
\DEELabstract{Although most of Brazil's electric energy comes from hydroelectric plants, 
thermal power plants are crucial for national energy security. However, with climate 
change caused by CO$_2$ emissions from various economic sectors, 
it is crucial to balance this security with CO$_2$ emissions in the Brazilian 
electricity sector. This work proposes estimating hourly CO$_2$ equivalent (CO$_2$-e) 
emissions from Brazilian thermal power plants using neural networks to support energy and 
environmental planning of the National Interconnected System. While the literature bases 
the modeling of these emissions on specific functions for each plant, emissions in Brazil are 
inventoried by Eletrobras and the Instituto de Energia e Meio Ambiente on an annual basis. 
To overcome this, this work discretizes the annual emissions of each plant into hourly estimates 
using nonlinear regression based on data provided by the Operador Nacional do Sistema Elétrico. 
This enables generalizing the estimates of hourly emissions from Brazilian thermal power plants. 
This work includes analyses for 39 thermal power plants with different generation capacities, 
operation cycles, and fuels, providing insights for energy, electric, and environmental planning 
in the country.} 

% Elementos pré-textuais
\ElementosPreTextuais
% Insere todos os elementos Pré-Textuais listados (capa, folha de rosto, ...)
    % capa                   % Obrigatório
    % folhaderosto           % Obrigatório
    % folhaaprovacao         % Obrigatório na versão final
	% dedicatoria            % Opcional
	% agradecimentos         % Opcional
    % epigrafe               % Opcional
    % resumo                 % Obrigatório
    % abstract               % Obrigatório

\newpage

%--------Lista de figuras--------
\ListaDeFiguras
\newpage
%--------Lista de tabelas--------
\ListaDeTabelas
\newpage
%--------Lista de Siglas e Abreviaturas--------
%----------Lista de Acronimos --------------------------------
\ifx\isenglish\undefined
\pretextualchapter{Lista de Siglas e Abreviaturas}
\else
\pretextualchapter{Acronyms}
\fi
%--------------------------

\begin{acronym}[itemsep=0pt]

    \acro{CO$_2$}{Dióxido de Carbono}
    \acro{CO$_2$e}{Dióxido de Carbono Equivalente}
    \acro{tCO$_2$e}{Toneladas de Dióxido de Carbono Equivalente}
    \acro{ONS}{Operador Nacional do Sistema Elétrico}
    \acro{IEMA}{Instituto de Energia e Meio Ambiente}
    \acro{SIN}{Sistema Interligado Nacional}
    \acro{UTE}{Usina Termelétrica}
    \acro{CEG}{Cadastro Específico de Geração}
    \acro{AMPL}{\emph{A Mathematical Programming Language}}
    \acro{IPOPT}{\emph{Interior Point Optimizer}}
    \acro{RNL}{Regressão Não Linear}
    \acro{MLP}{\emph{Multilayer Perceptron}}
    \acro{GRU}{\emph{Gated Recurrent Units }}
    \acro{LSTM}{\emph{Long Short-Term Memory}}
    \acro{ReLU}{\emph{Rectified Linear Unit}}
    \acro{MSE}{\emph{Mean Squared Error}}
    \acro{MAE}{\emph{Mean Absolute Error}}
    \acro{RMSE}{\emph{Root Mean Squared Error}}

\end{acronym}

%\addcontentsline{toc}{chapter}{Lista de Siglas e Abreviaturas}

\newpage
%--------Lista de Simbolos e Notação--------
\include{list_simbolos}
\newpage
%--------Sumário-------
\Sumario
\newpage

%%%%%%%%%%%%%%%%%%%%%%
%%Elementos Textuais%%
%%%%%%%%%%%%%%%%%%%%%%
\textual 

\chapter{Introdução}\label{chp:intro}

O setor elétrico brasileiro possui uma matriz predominantemente renovável, com destaque para as usinas hidrelétricas, que historicamente respondem pela maior parte da geração de energia elétrica no país. No entanto, períodos de estiagem prolongada, aliados ao crescimento da demanda e à necessidade de segurança energética, exigem o acionamento frequente de usinas termelétricas. Embora essas usinas sejam fundamentais para garantir o suprimento contínuo de energia, sua operação está associada à emissão de gases de efeito estufa, especialmente o \ac{CO$_2$}. As mudanças climáticas, impulsionadas pelas emissões de diversos setores econômicos, têm imposto desafios crescentes aos planejadores energéticos. No Brasil, o \ac{SIN} requer ferramentas que permitam equilibrar a segurança do suprimento elétrico com as metas de descarbonização assumidas pelo país em acordos internacionais. Nesse contexto, a estimação precisa das emissões de \ac{CO$_2$} das termelétricas em escala horária torna-se um insumo relevante para o despacho econômico-ambiental e para o planejamento da operação.

A literatura internacional apresenta modelos consolidados para a estimação de emissões horárias de usinas termelétricas, geralmente baseados em funções quadráticas ou exponenciais calibradas a partir de medições contínuas ou de dados de projeto de cada usina. No entanto, o cenário brasileiro apresenta uma particularidade: os dados de emissões de \ac{CO$_2$} das termelétricas são inventariados anualmente — e não horariamente — por instituições como o \ac{IEMA} e, anteriormente, pela Eletrobras. Essa disponibilidade exclusivamente anual de dados impõe um problema de natureza inversa: é necessário estimar o comportamento horário das emissões a partir de agregados anuais. Além disso, as termelétricas brasileiras operam com diferentes tipos de combustível (gás natural, óleo diesel, carvão mineral, biomassa), diferentes ciclos de operação (ciclo simples, ciclo combinado) e operam em patamares discretos de carga, com transições entre níveis de geração que afetam a eficiência e, consequentemente, as emissões. A literatura nacional ainda carece de metodologias sistemáticas para a discretização temporal das emissões anuais em estimativas horárias que considerem essas particularidades operacionais e que sejam generalizáveis para diferentes usinas.

Diante desse cenário, este trabalho propõe uma metodologia para estimar as emissões horárias de \ac{CO$_2$e} das usinas termelétricas brasileiras, combinando técnicas de otimização não-linear e redes neurais artificiais. A metodologia estrutura-se em duas etapas principais. Na primeira etapa, calibram-se os coeficientes de uma equação paramétrica de emissão (composta por termos quadrático e exponencial) por meio de um problema de minimização resolvido em \ac{AMPL} com o solver \ac{IPOPT}, utilizando como referência os totais anuais do \ac{IEMA}. Dois cenários são comparados: calibração padrão e calibração com regularização L2. Na segunda etapa, as emissões sintéticas horárias geradas tornam-se alvos para o treinamento de uma rede neural \ac{MLP}, que aprende diretamente a relação entre variáveis operacionais (geração, categoria de operação, fase de transição), características estáticas da usina (combustível, potência instalada, eficiência) e as emissões horárias. Mais especificamente, o trabalho desenvolve um algoritmo de classificação de categorias de geração e identificação de transições operacionais com base em limiares adaptativos, calibra os coeficientes da equação de emissão para 39 usinas termelétricas brasileiras no período de 2020 a 2024, treina e valida uma rede neural \ac{MLP} para estimar emissões horárias, e compara o desempenho dos cenários de calibração (padrão vs. L2) em termos de precisão e estabilidade.

Como contribuições principais, este trabalho oferece uma metodologia sistemática para discretização de emissões anuais em estimativas horárias aplicável a termelétricas com diferentes combustíveis e ciclos de operação; uma base de dados integrada combinando informações da \ac{ONS} e do \ac{IEMA} para 39 usinas em um período de 5 anos, com enriquecimento de \textit{features} temporais (categorias, transições, fases); uma comparação quantitativa entre calibração padrão e com regularização L2, fornecendo subsídios para a escolha da estratégia mais adequada em problemas de calibração inversa no setor elétrico; e um modelo de rede neural treinado para estimação horária de emissões, com análise de desempenho por categoria operacional e tipo de combustível. Essas contribuições podem apoiar o planejamento energético e ambiental do \ac{SIN}, fornecendo insumos horários para modelos de despacho econômico-ambiental e para a avaliação de políticas de descarbonização. % Texto do cap1.tex

\chapter{Revisão de Literatura}\label{chp:rev}

\section{Métricas de Emissões e \ac{CO$_2$e}}

A quantificação de emissões de gases de efeito estufa no setor energético brasileiro segue as diretrizes do \ac{IPCC}. Para todo gás \(g \in \{\text{CO}_2, \text{CH}_4, \text{N}_2\text{O}\}\) e combustível \(f\), as emissões são descritas em função da energia consumida e de fatores de emissão específicos: \(E_g = \sum_f \text{AD}_f \cdot \text{EF}_{f,g}\), onde \(\text{AD}_f\) são os dados de atividade (\textit{Activity Data}), correspondentes à energia consumida (em TJ), e \(\text{EF}_{f,g}\) é o fator de emissão (\textit{Emission Factor}) do gás \(g\) para o combustível \(f\) (em kg/TJ) \cite{ipcc2006}.

Após a quantificação das emissões individuais de cada gás, os inventários convertem o valor na unidade equivalente de dióxido de carbono (\ac{CO$_2$e}), utilizando os potenciais de aquecimento global (\textit{Global Warming Potential}) do \ac{AR5}:

\begin{equation}
\text{CO}_2e = \text{CO}_2 + 28 \cdot \text{CH}_4 + 265 \cdot \text{N}_2\text{O}
\label{eq:co2e}
\end{equation}

\cite{ipcc2014}. O \ac{IEMA} adota esta conversão em seus inventários anuais de emissões atmosféricas de usinas termelétricas, publicados a partir do ano-base 2020 \cite{iema2022, iema2023, iema2024}. Adicionalmente, o \ac{IEMA} divulga a taxa de emissão de cada usina em \ac{tCO$_2$e}/GWh, definida como a razão entre a emissão anual de \ac{CO$_2$e} e a energia elétrica gerada no ano, permitindo acompanhar o desempenho ambiental das plantas.

\section{Funções de Emissão Horária}

A estimativa de emissões de uma usina termelétrica em função da sua geração elétrica é um problema recorrente no despacho econômico-ambiental de sistemas de potência. A relação física subjacente decorre da cadeia de conversão energética: a queima do combustível libera energia térmica, que é convertida em trabalho mecânico e, posteriormente, em energia elétrica. Como o fator de emissão do combustível (kg CO\(_2\)/GJ) é aproximadamente constante para um dado tipo de combustível, a emissão apresenta a mesma forma funcional que o consumo de combustível \cite{ipcc2006}. Essa propriedade permite modelar diretamente a emissão como função da potência gerada, sem a necessidade de calcular o consumo intermediário de combustível.

A curva de entrada e saída (\textit{input-output curve}) de uma unidade térmica, que relaciona o aporte térmico \(H\) (em GJ/h) à potência elétrica gerada \(P\) (em MW), é classicamente representada por um polinômio de segundo grau \cite{wood2013power}:

\begin{equation}
H(P) = a + bP + cP^2
\label{eq:input_output}
\end{equation}

onde \(a\) representa o aporte térmico a vazio (\textit{no-load heat}), \(b\) o \textit{heat rate} incremental e \(c\) a curvatura associada à perda de eficiência em cargas elevadas. Multiplicando-se cada coeficiente pelo fator de emissão \(\rho\) (kg CO\(_2\)e/GJ), obtém-se diretamente a emissão horária como função da potência gerada:

\begin{equation}
E(P) = a + bP + cP^2 \quad [\text{tCO}_2\text{e}/\text{h}]
\label{eq:emissao_quadratica}
\end{equation}

onde \(a\), \(b\) e \(c\) são coeficientes de emissão intrínsecos a cada usina, determinados experimentalmente ou por ajuste a dados operacionais \cite{gent1971minimum, wood2013power}.

Essa formulação quadrática foi proposta por Gent e Lamont \cite{gent1971minimum} no contexto do despacho de mínima emissão de NOx e é amplamente adotada como modelo padrão para representar emissões em regime permanente \cite{kothari2011power}. No mesmo trabalho, os autores propuseram uma extensão com termo exponencial para capturar o aumento abrupto de emissões em altas cargas:

\begin{equation}
E(P) = a + bP + cP^2 + d \exp(eP)
\label{eq:emissao_exponencial}
\end{equation}

Essa formulação foi posteriormente consolidada na revisão de Talaq, El-Hawary e El-Hawary \cite{talaq1994summary}, que compilou os algoritmos de despacho econômico-ambiental desenvolvidos desde 1970. Para \ac{CO$_2$}, cuja emissão depende essencialmente da quantidade de combustível queimado e não da temperatura de combustão, a forma quadrática é geralmente suficiente. No entanto, a formulação com termo exponencial é adotada neste trabalho por sua maior flexibilidade para capturar não-linearidades observadas nos dados de diferentes usinas.

Além da forma polinomial, outras representações funcionais têm sido empregadas para capturar a degradação de eficiência em carga parcial. A forma logarítmica utiliza o fato de que o \textit{heat rate} específico cresce de maneira aproximadamente logarítmica conforme a carga se afasta do ponto de projeto \cite{eia2015}. A forma exponencial pura, \(E = a \exp(bP)\), penaliza de maneira mais acentuada operações em carga parcial reduzida, sendo empregada em estudos de flexibilidade operacional de usinas a carvão \cite{henderson2014}.

Cabe ressaltar que todas essas formulações pressupõem operação em regime permanente. Na prática, usinas termelétricas estão sujeitas a mudanças operacionais, como partidas (\textit{start-up}), paradas (\textit{shut-down}) e rampas de carga, que introduzem emissões adicionais não capturadas pelos modelos de regime estacionário \cite{kumar2012}. Dentre esses eventos, a partida é particularmente relevante sob a ótica ambiental: ao ligar uma caldeira, é necessário reaquecer a massa metálica do sistema, purgar gases residuais e estabilizar a combustão, consumindo combustível adicional enquanto a unidade ainda não atingiu a potência mínima para conexão ao sistema \cite{kumar2012}.

\section{Regressão Não Linear}

As formas funcionais apresentadas na seção anterior descrevem a emissão horária de uma termelétrica em função da potência gerada e de variáveis operacionais. Contudo, os coeficientes dessas funções são específicos de cada usina e precisam ser determinados a partir de dados operacionais. Quando a forma funcional contém termos exponenciais, logarítmicos ou racionais, os coeficientes aparecem de maneira não linear no modelo, e sua estimação requer técnicas de \ac{RNL} \cite{seber1989nonlinear}.

Em termos formais, dado um modelo paramétrico \(f(\mathbf{x}; \boldsymbol{\theta})\), onde \(\boldsymbol{\theta} \in \mathbb{R}^q\) é o vetor de parâmetros a ser estimado, o modelo é classificado como não linear nos parâmetros quando ao menos uma derivada parcial \(\partial f/\partial \theta_j\) depende do próprio vetor \(\boldsymbol{\theta}\). Nessa situação, não existe solução analítica fechada para a estimação dos coeficientes, sendo necessário recorrer a algoritmos iterativos de otimização numérica. Modelos lineares nos parâmetros, como polinômios (\(E = a + bP + cP^2\)), possuem solução analítica por Mínimos Quadrados Ordinários, mas por conveniência podem ser ajustados pelo mesmo algoritmo de otimização utilizado para os demais modelos \cite{seber1989nonlinear}.

A estimação dos parâmetros consiste em encontrar o vetor \(\hat{\boldsymbol{\theta}}\) que minimize uma função objetivo \(\mathcal{J}\), a qual quantifica a discrepância entre os valores observados e os preditos pelo modelo:

\begin{equation}
\hat{\boldsymbol{\theta}} = \arg \min_{\boldsymbol{\theta}} \mathcal{J}(\mathbf{y}, f(\mathbf{X}; \boldsymbol{\theta}))
\label{eq:estimacao_rnl}
\end{equation}

A escolha de \(\mathcal{J}\) depende das características do problema. A soma dos quadrados dos resíduos (\ac{MSE}) é a formulação clássica, mas alternativas como o \ac{MAE} e o \ac{MAPE} podem ser preferíveis em determinados cenários. O \ac{MSE} eleva os resíduos ao quadrado, amplificando o peso de valores discrepantes e puxando a solução para acomodá-los; o \ac{MAE} trata os desvios linearmente, resultando em ajuste mais robusto a \textit{outliers}. O \ac{MAPE}, por sua vez, normaliza cada resíduo pela magnitude da observação, sendo útil quando se deseja avaliar o erro em termos relativos \cite{myttenaere2016mean}.

No contexto de curvas de emissão de termelétricas, os coeficientes do modelo possuem significado físico e devem ser não negativos (\(\theta \geq 0\)), uma vez que representam taxas de emissão ou de consumo que são positivas por natureza \cite{wood2013power}.

Na prática, a convergência dos algoritmos iterativos utilizados em \ac{RNL} depende da escolha de valores iniciais e limites de busca adequados. A adoção de estimativas iniciais baseadas em conhecimento físico do problema assegura que a solução obtida possua significado físico e contribui para a aceleração da convergência \cite{seber1989nonlinear}.

\section{Geração de Dados Sintéticos}

Conforme apresentado na Introdução, os inventários do \ac{IEMA} fornecem emissões em resolução anual, enquanto os registros do \ac{ONS} estão disponíveis em base horária. Essa incompatibilidade temporal impede a aplicação direta de técnicas de regressão para modelar emissões horárias. A alternativa adotada neste trabalho consiste em gerar uma base de dados sintéticos horários de emissão, utilizando modelos paramétricos calibrados com os totais anuais disponíveis.

A geração de dados sintéticos busca criar dados que minimamente descrevam efeitos físicos e reais quando aqueles não estão disponíveis, seja por falta de medição ou por questões de confidencialidade. A ideia não é reproduzir cada observação individualmente, mas gerar amostras que preservem as relações estatísticas entre a variável resposta e os preditores presentes nos dados originais \cite{nowok2016synthpop}.

Na prática, isso requer a formulação de um modelo estatístico ou físico-matemático que represente o fenômeno de interesse. A partir desse modelo e de uma base de dados real, geram-se novas observações de modo que a base artificial resultante preserve as correlações e a variabilidade dos dados originais \cite{nowok2016synthpop}.

Assim, a criação de dados sintéticos, quando bem sucedida, permite realizar testes e experimentos sobre uma base de dados que, se tratada apenas com os valores originais, seria insuficiente para estudo. Porém, como essa geração artificial depende diretamente do modelo utilizado e dos dados originais, quanto melhor o modelo descrever a realidade e quanto menor o ruído nos dados de entrada, maior será a confiabilidade dos dados sintéticos. Ainda assim, os dados sintéticos gerados podem reproduzir o formato geral das distribuições, distorcer dependências ou não preservar valores extremos e eventos esparsos \cite{nowok2016synthpop, stenger2024synthetic}.

A geração de bases sintéticas para variáveis dependentes do tempo ou altamente correlacionadas apresenta desafios adicionais. Uma série temporal sintética pode, por exemplo, reproduzir a distribuição marginal de cada variável, mas falhar em replicar padrões de autocorrelação ou relações de defasagem \cite{stenger2024synthetic}.

Em síntese, dados sintéticos constituem um recurso valioso para experimentação e validação quando dados reais são insuficientes, desde que interpretados com cautela — embora sejam capazes de aproximar a estrutura estatística do conjunto original, podem introduzir viés ou enfraquecer relações dinâmicas entre variáveis. A documentação completa do método de geração, das escolhas de parâmetros e de suas limitações é, portanto, essencial para o uso confiável desses dados \cite{nowok2016synthpop, stenger2024synthetic}.

\section{Métricas de Avaliação}

Para quantificar a qualidade do ajuste e a capacidade preditiva dos modelos, são empregadas métricas estatísticas complementares. Essas métricas são utilizadas tanto na etapa de ajuste dos modelos de \ac{RNL} (como função objetivo a ser minimizada) quanto na avaliação da rede neural. Em todas as definições a seguir, \(y_i\) denota o valor observado, \(\hat{y}_i\) o valor predito pelo modelo, \(\bar{y}\) a média das observações e \(N\) o número total de amostras.

\subsection{Erro Absoluto Médio (\ac{MAE})}

O \ac{MAE} mensura a média aritmética dos desvios absolutos entre valores preditos e observados. Por tratar os desvios de forma linear, é mais robusto a \textit{outliers} do que o \ac{MSE} \cite{willmott2005advantages}:

\begin{equation}
\text{MAE} = \frac{1}{N} \sum_{i=1}^{N} |y_i - \hat{y}_i|
\label{eq:mae}
\end{equation}

\subsection{Erro Quadrático Médio (\ac{MSE})}

O \ac{MSE} calcula a média dos quadrados dos desvios. Por elevar os erros ao quadrado, penaliza mais severamente erros de grande magnitude \cite{bickel2015mathematical}:

\begin{equation}
\text{MSE} = \frac{1}{N} \sum_{i=1}^{N} (y_i - \hat{y}_i)^2
\label{eq:mse}
\end{equation}

\subsection{Raiz do Erro Quadrático Médio (\ac{RMSE})}

O \ac{RMSE} é obtido pela raiz quadrada do \ac{MSE}, apresentando o erro na mesma unidade dimensional da variável resposta, o que facilita a interpretação física do desvio \cite{chai2014root}:

\begin{equation}
\text{RMSE} = \sqrt{\text{MSE}}
\label{eq:rmse}
\end{equation}

\subsection{Erro Percentual Absoluto Médio (\ac{MAPE})}

O \ac{MAPE} expressa o desvio em termos relativos (porcentagem), sendo útil para comparar o desempenho do modelo em diferentes escalas de magnitude \cite{myttenaere2016mean}:

\begin{equation}
\text{MAPE} = \frac{100}{N} \sum_{i=1}^{N} \left| \frac{y_i - \hat{y}_i}{y_i} \right|
\label{eq:mape}
\end{equation}

\subsection{Razão de Estabilidade}

A razão \(R = \text{MSE} / \text{MAE}^2\) é empregada como indicador da estabilidade do modelo. Valores baixos de \(R\) indicam que os erros são consistentes, enquanto valores altos sugerem que erros grandes ocorrem com frequência não desprezível:

\begin{equation}
R = \frac{\text{MSE}}{\text{MAE}^2}
\label{eq:estabilidade}
\end{equation}

A interpretação adotada neste trabalho é:
\begin{itemize}
    \item \(R < 3\): modelo estável, erros concentrados em torno da média;
    \item \(3 \leq R \leq 5\): instabilidade moderada, presença de alguns \textit{outliers};
    \item \(R > 5\): presença significativa de \textit{outliers}, modelo pouco confiável.
\end{itemize} % Texto do cap2.tex

\chapter{Metodologia}\label{chp:met}

\section{Visão Geral}

A estimação horária de emissões de \ac{CO$_2$} em usinas termelétricas configura-se como um problema inverso: dispõe-se de medições anuais agregadas, fornecidas pelos inventários do \ac{IEMA}, mas necessita-se do comportamento horário para aplicações em planejamento energético e despacho econômico-ambiental. Além disso, a relação entre geração horária e emissão é intrinsecamente não-linear \cite{gent1971minimum, talaq1994summary}.

Para enfrentar esses desafios, a metodologia proposta estrutura-se em duas etapas complementares. Na primeira etapa, resolve-se um problema de otimização não-linear para calibrar os coeficientes de uma equação paramétrica de emissão horária, utilizando como referência os totais anuais disponibilizados pelo \ac{IEMA}. Para mitigar o mal condicionamento desse problema inverso, aplica-se regularização L2 (Ridge), sendo comparados dois cenários: calibração padrão e calibração com regularização.

Na segunda etapa, as emissões sintéticas horárias geradas pela equação calibrada tornam-se o alvo (\textit{target}) para o treinamento de uma rede neural \ac{MLP}. A rede aprende diretamente a relação entre as variáveis operacionais (geração, categoria de operação, fase de transição), as características estáticas da usina (combustível, potência instalada, eficiência) e a emissão horária, sem incorporar explicitamente a equação física na função de perda.

Essa abordagem em cascata permite: (i) garantir consistência com os totais anuais do \ac{IEMA}; (ii) capturar não-linearidades que a equação paramétrica pode não expressar plenamente; e (iii) avaliar, por meio da comparação entre os cenários padrão e L2, qual estratégia de calibração produz alvos sintéticos mais coerentes para o aprendizado supervisionado.

\section{Fontes de Dados e Descrição das Variáveis}

Foram empregadas duas fontes principais de dados, abrangendo o período de 2020 a 2024.

\subsection{Operador Nacional do Sistema Elétrico (\ac{ONS})}

Por meio da iniciativa Dados Abertos, obtém-se a geração horária de todas as usinas do \ac{SIN}. Os registros incluem: identificação da usina, data, hora e geração horária em MW. O conjunto é filtrado para reter apenas usinas termelétricas \cite{ons_dados}.

\subsection{Instituto de Energia e Meio Ambiente (\ac{IEMA})}

Os inventários anuais do \ac{IEMA} fornecem \cite{iema2022, iema2023, iema2024}:
\begin{itemize}
    \item Emissões totais anuais de \ac{tCO$_2$e} por usina;
    \item Características cadastrais fixas: potência instalada (MW), fator de capacidade (\%), eficiência energética (\%), tipo de combustível e ciclo de operação (ciclo combinado, ciclo simples, cogeração).
\end{itemize}

Usinas de cogeração são excluídas da análise por não apresentarem, nos inventários do \ac{IEMA}, informações completas de potência instalada e eficiência — características indispensáveis à modelagem. Após a integração das fontes e a remoção de registros faltantes, restam 39 usinas termelétricas com dados completos para o período 2020-2024. As emissões anuais são mantidas em \ac{tCO$_2$e}, e a geração horária em MW.

\section{Tratamento das Sazonalidades Temporais}

As variáveis temporais \textit{Índice} (posição da hora dentro do ano) e \textit{Ano} não podem ser tratadas como grandezas ordinais lineares, pois a relação entre instantes próximos ao final e ao início de um ano ou período não é capturada por uma codificação escalar convencional.

Para preservar a periodicidade natural do calendário, adota-se uma codificação polar (seno/cosseno). Dado um índice \(i \in [1, H]\), onde \(H \in \{8760, 8784\}\) representa as horas do respectivo ano (considerando anos bissextos), definem-se as seguintes transformações:

\begin{equation}
\text{Seno\_Índice}(i) = \sin\left(\frac{2\pi i}{H}\right), \qquad
\text{Cosseno\_Índice}(i) = \cos\left(\frac{2\pi i}{H}\right)
\label{eq:codificacao_indice}
\end{equation}

Analogamente, para o ano \(a \in [2020, 2024]\), define-se:

\begin{equation}
\text{Seno\_Ano}(a) = \sin\left(\frac{2\pi (a - 2020)}{4}\right), \qquad
\text{Cosseno\_Ano}(a) = \cos\left(\frac{2\pi (a - 2020)}{4}\right)
\label{eq:codificacao_ano}
\end{equation}

Essa representação mapeia valores temporais para pontos sobre o círculo unitário, de modo que instantes próximos (ex: hora 8760 e hora 1) tornam-se vizinhos no espaço polar, refletindo sua continuidade temporal.

\section{Modelagem das Transições Operacionais}

Termelétricas operam tipicamente em patamares discretos de carga, não em regime continuamente variável. Para capturar esse comportamento, define-se um sistema de categorias de geração com base na proximidade relativa entre valores de potência.

\subsection{Definição de Categorias}

Dada a série de geração horária \(G = [g_1, g_2, \ldots, g_T]\) em MW, valores inferiores a 20 MW são classificados na categoria 0 (carga muito baixa ou desligamento). Para os demais valores, dois pontos pertencem à mesma categoria se sua diferença relativa for inferior ou igual a um limiar \(\theta = 15\%\) do maior valor do grupo.

Formalmente, \(g_i\) e \(g_j\) pertencem à mesma categoria se:

\begin{equation}
\frac{|g_i - g_j|}{\max(g_i, g_j)} \leq \theta
\label{eq:diferenca_relativa}
\end{equation}

O limiar de 15\% foi definido com base na análise exploratória da distribuição dos incrementos de carga nos dados do \ac{ONS}, representando o valor típico que separa regimes operacionais distintos sem fragmentar artificialmente a operação.

\subsection{Identificação de Transições}

Seja \(c_t\) a categoria no instante \(t\). Define-se a moda \(m_t\) das categorias em uma janela de três horas anteriores e três horas posteriores, excluindo o próprio instante:

\begin{equation}
m_t = \text{moda}\{c_{t-3}, c_{t-2}, c_{t-1}, c_{t+1}, c_{t+2}, c_{t+3}\}
\label{eq:transicao}
\end{equation}

O instante \(t\) é classificado como \textbf{transição} se \(c_t \neq m_t\). Este critério captura períodos de rampa de carga (partidas, paradas e variações controladas), nos quais o fator de emissão horária pode diferir substancialmente do regime estabilizado \cite{kumar2012}.

\subsection{Duração e Fase da Transição}

Para cada evento de transição contínuo (sequência de horas consecutivas classificadas como transição), definem-se:
\begin{itemize}
    \item \textbf{Duração} \(d\): número de horas consecutivas do evento;
    \item \textbf{Fase} \(f\): posição relativa dentro do evento, classificada como 1 (início), 2 (meio) ou 3 (fim). Transições com duração \(d = 1\) recebem \(f = 1\) (início e fim coincidentes).
\end{itemize}

Essas variáveis permitem que o modelo compreenda a transição como um processo contínuo, com dinâmica de emissão potencialmente distinta ao longo de suas fases.

\section{Formulação do Problema de Calibração}

A relação entre geração horária e emissão horária é modelada por uma equação com termo quadrático e termo exponencial, conforme proposta na literatura \cite{gent1971minimum, talaq1994summary}:

\begin{equation}
E_k^{\text{horária}}(P_{G_i}) = \alpha_k P_{G_i}^2 + \beta_k P_{G_i} + \gamma_k + \omega_k e^{\mu_k P_{G_i}}
\label{eq:emissao_horaria}
\end{equation}

onde:
\begin{itemize}
    \item \(P_{G_i}\) representa a geração em MW, dividida por 100 para normalização dos dados, garantindo estabilidade numérica no termo exponencial;
    \item \(\alpha_k, \beta_k, \gamma_k, \omega_k, \mu_k\) são coeficientes a serem calibrados por usina;
    \item O subscrito \(k\) indica a usina.
\end{itemize}

\subsection{Calibração Padrão}

Os coeficientes são obtidos resolvendo o seguinte problema de minimização, que busca minimizar o erro quadrático médio entre a emissão anual estimada (somatório das emissões horárias) e a emissão anual real fornecida pelo \ac{IEMA}:

\begin{equation}
\min_{\alpha,\beta,\gamma,\omega,\mu} \frac{1}{N} \sum_{j=1}^{N} \left( \sum_{i=1}^{H} E_{\text{usina}}^{i}(P_{G_i}) - E_{\text{usina}}^{j} \right)^2
\label{eq:minimizacao_padrao}
\end{equation}

sujeito às restrições físicas:

\begin{equation}
\alpha \geq 0, \quad \mu \leq \mu_{\max}
\label{eq:restricoes_fisicas}
\end{equation}

onde:
\begin{itemize}
    \item \(N\) é o número de usinas (\(N = 39\));
    \item \(H\) é o número de horas do respectivo ano (\(H = 8760\) para anos não bissextos, \(H = 8784\) para anos bissextos);
    \item \(E_{\text{usina}}^{j}\) é a emissão anual real do \ac{IEMA} para a usina \(j\);
    \item \(\mu_{\max}\) é um limite superior para evitar \textit{overflow} numérico do termo exponencial (definido empiricamente em \(0,5\)).
\end{itemize}

Apenas o coeficiente \(\alpha\) foi explicitamente restringido à não-negatividade, por representar o termo quadrático associado ao crescimento da emissão em cargas elevadas. Os coeficientes \(\gamma\) e \(\omega\) foram mantidos livres para permitir maior flexibilidade durante a calibração do problema inverso \cite{wood2013power}. O coeficiente \(\mu\) recebe um limite superior \(\mu_{\max} = 0,5\) para evitar crescimento excessivo do termo exponencial, que poderia causar instabilidade numérica.

\subsection{Regularização L2 (Ridge)}

O problema de calibração inversa é intrinsecamente mal condicionado, podendo produzir coeficientes excessivamente elevados que se cancelam mutuamente (\textit{overfitting} na calibração). Para mitigar esse efeito, introduz-se um termo de regularização L2 (Tikhonov) à função objetivo \cite{tikhonov1977solutions}:

\begin{equation}
\min_{\alpha,\beta,\gamma,\omega,\mu} \frac{1}{N} \sum_{j=1}^{N} \left( \sum_{i=1}^{H} E_{\text{usina}}^{i}(P_{G_i}) - E_{\text{usina}}^{j} \right)^2 + \lambda \left( \alpha^2 + \beta^2 + \gamma^2 + \omega^2 + \mu^2 \right)
\label{eq:regularizacao_l2}
\end{equation}

sujeito às mesmas restrições da calibração padrão (Equação \ref{eq:restricoes_fisicas}).

O parâmetro \(\lambda \geq 0\) controla a intensidade da regularização: \(\lambda = 0\) recupera o problema padrão (Equação \ref{eq:minimizacao_padrao}), enquanto valores positivos penalizam coeficientes de grande magnitude, promovendo soluções mais estáveis e generalizáveis. O valor de \(\lambda\) deve ser selecionado de forma a equilibrar a redução do erro de calibração e a estabilidade dos coeficientes. Neste trabalho, adotou-se \(\lambda = 30\), conforme justificado na Seção 4.3.3.

\section{Emissões Sintéticas como Alvos para Aprendizado}

Uma vez calibrados os coeficientes (para os cenários padrão e com regularização L2), aplica-se a Equação \ref{eq:emissao_horaria} sobre a série horária de geração para produzir as \textbf{emissões sintéticas horárias} \(\hat{E}_t\). Essas emissões sintéticas tornam-se a variável alvo para o treinamento supervisionado da rede neural.

É importante observar que a rede neural, ao ser treinada com esses alvos sintéticos, herda os erros de modelagem da equação paramétrica e da calibração. Conforme discutido na literatura \cite{nowok2016synthpop, stenger2024synthetic}, os dados sintéticos gerados por modelos paramétricos refletem integralmente as premissas e limitações do modelo utilizado. Portanto, a comparação entre os cenários padrão e L2 não avalia apenas a qualidade da rede, mas também qual estratégia de calibração produz alvos mais consistentes para o aprendizado supervisionado.

A fidelidade da base sintética depende diretamente da qualidade do ajuste e da capacidade do modelo paramétrico em representar o comportamento real de emissão da usina.

\section{Arquitetura da Rede Neural}

Optou-se por uma arquitetura do tipo \ac{MLP} em detrimento de redes recorrentes (\ac{LSTM}, \ac{GRU}) porque a emissão horária, no modelo proposto, depende essencialmente da geração no mesmo instante \(t\), e não de uma sequência temporal longa. A \ac{MLP} é suficiente, mais simples e menos suscetível a \textit{overfitting} com o volume de dados disponível \cite{goodfellow2016deep}.


\subsection{Modelos Individualizados por Usina}

Um aspecto fundamental da metodologia é que um modelo de rede neural \ac{MLP} separado é treinado para cada uma das 39 usinas. Esta abordagem individualizada é necessária porque cada usina possui características operacionais próprias: diferentes tipos de combustível, eficiências energéticas, potências instaladas e perfis de transição entre patamares de carga. Um modelo único para todas as usinas não seria capaz capturar essas particularidades. Portanto, o processo descrito na Seção 3.7.1 é repetido individualmente para cada usina, resultando em 39 modelos de rede neural distintos ao final do treinamento.

\subsection{Estrutura das Camadas}

A Tabela \ref{tab:arquitetura} apresenta a arquitetura adotada.

\begin{table}[H]
    \centering
    \renewcommand{\arraystretch}{1.5}
    \begin{tabular}{|c|c|c|c|}
        \hline
        \textbf{Camada} & \textbf{Neurônios} & \textbf{Ativação} & \textbf{Regularização} \\
        \hline
        Entrada & \(d_{\text{in}}\) & -- & -- \\
        Oculta 1 & 128 & \ac{ReLU} & Dropout (0,3) \\
        Oculta 2 & 64 & \ac{ReLU} & Dropout (0,3) \\
        Oculta 3 & 32 & \ac{ReLU} & Dropout (0,3) \\
        Saída & 1 & \ac{ReLU} & -- \\
        \hline
    \end{tabular}
    \caption{Arquitetura da Rede Neural MLP}
    \label{tab:arquitetura}
\end{table}

A arquitetura em funil (128 $\rightarrow$ 64 $\rightarrow$ 32) reduz progressivamente a capacidade representacional, atuando como uma forma implícita de regularização e prevenindo \textit{overfitting}. A função de ativação \ac{ReLU} é escolhida por sua simplicidade computacional e por mitigar o problema do gradiente evanescente em comparação com funções sigmoidais \cite{goodfellow2016deep}. O \textit{dropout} com taxa de \(0,3\) desativa aleatoriamente \(30\%\) dos neurônios a cada iteração de treinamento, forçando a rede a aprender representações redundantes e robustas \cite{srivastava2014dropout}.

\subsection{Estratégias de Regularização Adicionais}

Além do \textit{dropout}, empregam-se as seguintes estratégias de regularização:

\begin{itemize}
    \item \textbf{Early stopping}: o treinamento é interrompido quando a perda no conjunto de validação não apresenta melhora por um número pré-definido de épocas (paciência), restaurando os pesos da melhor época. Isso evita \textit{overfitting} ao interromper o treinamento antes que a rede memorize ruídos \cite{goodfellow2016deep}.

    \item \textbf{Gradient clipping}: a norma L2 do vetor de gradientes é limitada a um valor máximo antes da atualização dos pesos. Esta técnica previne a explosão dos gradientes, fenômeno que pode ocorrer quando a função de perda apresenta regiões de alta curvatura ou quando amostras com pesos elevados produzem gradientes desproporcionais \cite{goodfellow2016deep}.
\end{itemize}

\section{Variáveis de Entrada e Pré-processamento Conceitual}

As variáveis de entrada da rede neural são agrupadas em duas categorias: estáticas e dinâmicas.

\subsection{Variáveis Estáticas}

Variáveis estáticas são invariantes no tempo para cada usina:
\begin{itemize}
    \item Tipo de combustível (categórica, sem ordem intrínseca);
    \item Ciclo de operação (categórica);
    \item Potência instalada (MW) — grandeza com significado físico direto;
    \item Eficiência energética (\%) — grandeza com significado físico direto;
    \item Fator de capacidade (\%) — grandeza com significado físico direto.
\end{itemize}

\subsection{Variáveis Dinâmicas}

Variáveis dinâmicas variam a cada hora:
\begin{itemize}
    \item Geração horária (MW);
    \item Categoria de geração (categórica, proveniente da classificação descrita na Seção 4);
    \item Duração da transição (h);
    \item Fase da transição (1, 2 ou 3);
    \item Índice temporal (codificado polarmente conforme Seção 3);
    \item Ano (codificado polarmente conforme Seção 3).
\end{itemize}

\subsection{Tratamento das Variáveis}

Variáveis categóricas (combustível, ciclo de operação, categoria de geração) são transformadas por \textit{one-hot encoding}, que as converte em vetores binários sem impor ordenação artificial.

Variáveis contínuas de diferentes ordens de grandeza (geração, duração, fase) requerem padronização para que todas contribuam igualmente durante o treinamento. Adota-se a padronização \textit{z-score}:

\begin{equation}
x_{\text{pad}} = \frac{x - \mu}{\sigma}
\label{eq:padronizacao}
\end{equation}

onde \(\mu\) e \(\sigma\) são a média e o desvio padrão estimados \textbf{exclusivamente} nos dados de treinamento, para evitar vazamento de informações do futuro.

\section{Métricas de Avaliação}

Para avaliar a qualidade das estimativas, empregam-se as seguintes métricas \cite{chai2014root, willmott2005advantages, myttenaere2016mean}.

\subsection{Erro Absoluto Médio (\ac{MAE})}

\begin{equation}
\text{MAE} = \frac{1}{n} \sum_{i=1}^{n} |y_i - \hat{y}_i|
\label{eq:mae}
\end{equation}
Interpretável diretamente na unidade original (tCO\(_2\)/h).

\subsection{Erro Quadrático Médio (\ac{MSE})}

\begin{equation}
\text{MSE} = \frac{1}{n} \sum_{i=1}^{n} (y_i - \hat{y}_i)^2
\label{eq:mse}
\end{equation}
Penaliza mais severamente erros grandes, sendo útil para detectar a presença de \textit{outliers}.

\subsection{Raiz do Erro Quadrático Médio (\ac{RMSE})}

\begin{equation}
\text{RMSE} = \sqrt{\text{MSE}}
\label{eq:rmse}
\end{equation}
Devolve a métrica à unidade original (tCO\(_2\)/h), combinando a sensibilidade a \textit{outliers} do \ac{MSE} com a interpretabilidade do \ac{MAE}.

\subsection{Razão de Estabilidade}

A razão $R = \text{MSE} / \text{MAE}^2$ é empregada como indicador da estabilidade do modelo. Esta razão está diretamente relacionada à distribuição dos erros e à presença de \textit{outliers}. Para uma distribuição normal de erros, a razão converge para $\pi/2 \approx 1,57$, pois $\text{MSE} = \sigma^2$ e $\text{MAE} = \sigma \sqrt{2/\pi}$, resultando em $\text{MSE}/\text{MAE}^2 = \pi/2$.

Com base nesta propriedade e em considerações práticas, adotam-se os seguintes limiares:

\begin{itemize}
    \item $R < 3$: modelo estável, erros concentrados em torno da média (até aproximadamente o dobro do valor teórico);
    \item $3 \leq R \leq 5$: instabilidade moderada, presença de alguns \textit{outliers};
    \item $R > 5$: presença significativa de \textit{outliers}, modelo pouco confiável.
\end{itemize}

Estes limiares foram definidos com base na análise exploratória dos dados e na experiência prática do autor, sendo utilizados como referência para comparação entre os cenários avaliados. % Texto do cap3.tex

\chapter{Desenvolvimento}\label{chp:des}

\section{Organização Geral da Implementação}

Todo o código foi desenvolvido na linguagem Python, versão 3.10.12, utilizando as bibliotecas Pandas e NumPy para manipulação dos dados, \ac{AMPL} (via biblioteca \emph{amplpy}) para a calibração dos coeficientes, e PyTorch para a construção e treinamento da rede neural.

O desenvolvimento foi estruturado em três módulos principais:

\begin{itemize}
    \item \textbf{Módulo 1 – Pré-processamento:} filtragem das usinas, organização temporal, classificação de categorias, identificação de transições e cálculo de duração/fase.

    \item \textbf{Módulo 2 – Calibração (AMPL/IPOPT):} modelagem do problema de otimização, integração Python-AMPL, ajuste do parâmetro de regularização e geração das emissões sintéticas.

    \item \textbf{Módulo 3 – Rede Neural (PyTorch):} construção da arquitetura MLP, ponderação das amostras, treinamento com early stopping e gradient clipping.
\end{itemize}

Todos os experimentos foram executados em um computador com processador AMD Ryzen 5 5600X (12 núcleos), 16 GB de RAM e GPU NVIDIA RTX 5070 (12 GB), sob sistema operacional Windows 11.

\section{Implementação do Pré-Processamento}

\subsection{Filtragem das Usinas e Organização Temporal}

O primeiro passo consistiu em integrar as duas fontes de dados e garantir a consistência do conjunto para o período analisado (2020 a 2024). A partir dos dados da \ac{ONS}, filtrou-se inicialmente apenas as usinas termelétricas. 

Em seguida, implementou-se um algoritmo para reter apenas as usinas que possuíam registros completos em todos os inventários anuais do \ac{IEMA} para o período de interesse. Usinas de cogeração foram automaticamente excluídas devido à ausência de informações de potência instalada e eficiência energética nos inventários.

Após a filtragem, cada usina foi identificada unicamente por seu \ac{CEG}. Os dados horários de geração de cada usina foram então separados por ano e concatenados em uma única série temporal por usina, preservando duas colunas fundamentais: ``Ano'' (2020 a 2024) e ``Índice'' (posição da hora dentro do ano, variando de 1 a 8760, ou 8784 para anos bissextos). Esta organização permitiu manter a estrutura temporal necessária para etapas subsequentes.

\subsection{Classificação das Categorias de Geração}

O algoritmo de classificação de categorias foi implementado conforme o procedimento descrito na Metodologia. A implementação percorre todos os valores de geração acima de 20 MW, agrupando aqueles cuja diferença relativa em relação à média do grupo não excede o limiar de 15\%. Após a formação inicial dos grupos, realiza-se uma etapa iterativa de mesclagem, na qual grupos cujas médias diferem em menos de 15\% são combinados. Por fim, valores abaixo do limiar de 20 MW são atribuídos à categoria 1, pois representam carga muito baixa ou desligamento.

\subsection{Identificação de Transições}

Para identificar horas em transição, implementou-se uma função que percorre o vetor de categorias com uma janela deslizante de três horas anteriores e três horas posteriores. Para cada hora central da janela, calcula-se a moda das categorias nos seis instantes vizinhos. Se a categoria da hora central for diferente da moda dos vizinhos, a hora é marcada como transição.

As horas nas bordas do ano utilizam apenas as informações disponíveis para determinação de sua classificação.

\section{Implementação da Calibração no \ac{AMPL}}

\subsection{Estrutura do Modelo de Otimização}

O problema de minimização foi modelado na linguagem \ac{AMPL}. Foram criados dois arquivos: um arquivo de modelo (\emph{.mod}) contendo as variáveis, a função objetivo e as restrições; e um arquivo de comandos (\emph{.run}) especificando as opções do solver \ac{IPOPT}.

A função objetivo implementa a Equação \textbf{INSERIR NUMERAÇÃO} para o cenário padrão, com a adição do termo de regularização L2. Os dados de geração horária e a emissão anual alvo são fornecidos separadamente.

\subsection{Integração Python-\ac{AMPL}}

Para automatizar a calibração das 39 usinas, utilizou-se a biblioteca \emph{amplpy}. O fluxo de integração foi o seguinte:

\begin{itemize}
    \item Para cada usina, extraem-se os dados de geração horária e a emissão anual correspondente do \ac{IEMA};

    \item Os dados são escritos em um arquivo \emph{.dat} para cada usina;

    \item O modelo \ac{AMPL} é iterado em cada usina, com seu respectivo arquivo \emph{.dat} importado;

    \item O solver \ac{IPOPT} é executado;

    \item Os coeficientes otimizados são extraídos e armazenados em um dicionário Python.
\end{itemize}

Este processo foi executado em paralelo utilizando a biblioteca \emph{multiprocessing}, reduzindo o tempo total de calibração de aproximadamente 8 horas em processamento serial para cerca de 2 horas.

\subsection{Ajuste do Parâmetro de Regularização $\lambda$}

O parâmetro $\lambda$ foi ajustado por meio de testes empíricos com os valores 1, 10, 30 e 100. Para cada valor, calibrou-se um subconjunto de 10 usinas representativas e avaliou-se o erro \ac{MAE} médio e a variância dos coeficientes entre anos consecutivos.

O valor $\lambda = 30$ foi selecionado por apresentar o melhor equilíbrio. Valores superiores como $\lambda=100$ introduziram viés excessivo, enquanto valores inferiores como $\lambda=1$ não foram suficientes para estabilizar os coeficientes.

\subsection{Geração das Emissões Sintéticas Horárias}

Uma vez calibrados os coeficientes para cada usina, aplicou-se a Equação \textbf{INSERIR NUMERO} sobre a série horária de geração para produzir as emissões sintéticas. Este procedimento foi executado separadamente para os dois cenários, padrão e com regularização L2.

\section{Preparação das Features para a Rede Neural}

\subsection{Tratamento das Variáveis}

Antes do treinamento, as variáveis de entrada foram preparadas conforme as Tabelas \ref{tab:var_est} e \ref{tab:var_din}.

\begin{table}[H]
    \center
    \renewcommand{\arraystretch}{1.5}
    \begin{tabular}{|c|c|c|}
        \hline
        \textbf{Variável}     & \textbf{Tratamento} & \textbf{Justificativa}                 \\ \hline
        Combustível           & One-Hot Encoding    & Categoria nominal sem ordem intrínseca \\ \hline
        Ciclo de operação     & One-Hot Encoding    & Categoria nominal                      \\ \hline
        Potência instalada    & Sem normalização    & Possui significado físico direto       \\ \hline
        Eficiência energética & Sem normalização    & Possui significado físico direto       \\ \hline
        Fator de capacidade   & Sem normalização    & Possui significado físico direto       \\ \hline
    \end{tabular}
    \caption{Tratamento das Variáveis Estáticas}
    \label{tab:var_est}
\end{table}

\begin{table}[H]
    \center
    \renewcommand{\arraystretch}{1.5}
    \begin{tabularx}{\textwidth}{|c|c|X|}
        \hline
        \textbf{Variável}        & \textbf{Tratamento}              & \textbf{Justificativa}                          \\ \hline
        Geração (MW)             & Standard Scaler                  & Ajustado após divisão treino/validação/teste                       \\ \hline
        Duração da transição (h) & Standard Scaler                  & Ajustado para rede neural                       \\ \hline
        Fase da transição        & Standard Scaler                  & Ajustado para rede neural                       \\ \hline
        Índice temporal          & Seno/Cosseno (Codificação polar) & Preserva periodicidade anual                    \\ \hline
        Ano                      & Seno/Cosseno (Codificação polar) & Permite capturar tendências anuais não lineares \\ \hline
        \caption{Tratamento das Variáveis Dinâmicas}\label{tab:var_din}
    \end{tabularx}
\end{table}

\subsection{Cálculo da Duração e Fase da Transição}

Ademais, o cálculo da duração e fase das gerações de transição foi implementado varrendo o vetor binário de transição e identificando blocos consecutivos onde o valor é 1. Para cada bloco, a duração é o comprimento do bloco, e a fase é atribuída conforme a posição relativa: primeira hora = 1, última hora = 3, horas intermediárias = 2. Transições de duração 1 recebem fase = 1.

\section{Implementação da Rede Neural}

\subsection{Construção da Arquitetura \ac{MLP}}

A rede neural foi implementada utilizando PyTorch. A escolha desta biblioteca deveu-se à possibilidade de executar o treinamento, validação e teste utilizando a GPU da máquina, acelerando significativamente o processamento. 

A arquitetura seguiu a estrutura definida na Tabela \textbf{INSERIR NUMERO}: três camadas ocultas com 128, 64 e 32 neurônios, ativação \ac{ReLU}, \emph{dropout} de 0,3 após cada camada oculta, e camada de saída com um neurônio e ativação \ac{ReLU} para garantir não-negatividade das emissões. O horizonte de previsão foi definido como 1 hora.

\subsection{Ponderação Diferenciada para Transições}

Para melhorar a capacidade do modelo em aprender eventos de transição, implementou-se um sistema de ponderação na função de perda. Cada amostra recebeu um peso individual: 1,0 para horas em regime estável e 5,0 para horas classificadas como transição.

O valor 5 foi escolhido empiricamente: pesos menores como 2 ou 3 não produziram melhora significativa no erro das transições, enquanto pesos maiores como 10 causaram overfitting a esses eventos.

\subsection{Early Stopping e Gradient Clipping}

O método early stopping foi implementado manualmente no loop de treinamento, visando obter os melhores valores para os pesos. Assim, se nenhuma melhora fosse observada após 50 épocas consecutivas, o treinamento é interrompido e os pesos da melhor época eram restaurados. 

O gradient clipping atua juntamente, calculando o gradiente e atualizando os pesos para mantê-los no limiar de [-1, 1], evitando que os gradientes se tornem excessivamente grandes durante a retropropagação.

\section{Divisão dos Dados e Configuração do Treinamento}

A divisão temporal dos dados foi configurada conforme a Tabela \ref{tab:div_dados}

\begin{table}[H]
    \center
    \renewcommand{\arraystretch}{1.5}
    \begin{tabular}{|c|c|c|}
        \hline
        \textbf{Conjunto} & \textbf{Período} & \textbf{Finalidade}                        \\ \hline
        Treinamento       & 2020, 2021, 2022 & Ajuste dos pesos                           \\ \hline
        Validação         & 2023             & Early stopping, seleção de hiperparâmetros \\ \hline
        Teste             & 2024               & Avaliação final do modelo                  \\ \hline
    \end{tabular}
    \caption{Divisão Temporal dos Dados}
    \label{tab:div_dados}
\end{table}

Os hiperparâmetros de treinamento foram configurados conforme a Tabela \ref{tab:hiperparam}

\begin{table}[H]
    \centering
    \renewcommand{\arraystretch}{2}
    \begin{tabular}{|c|c|c|}
    \hline
        \textbf{Métrica}       & \textbf{Fórmula}           & \textbf{Justificativa}                             \\ \hline
        MAE (tCO$_2$/h)                    & $\frac{1}{n}\sum y_i - \hat{y}_i$ & Interpretação direta na unidade original  \\ \hline
        MSE (tCO$_2$/h)$^2$ & $\frac{1}{n}\sum (y_i - \hat{y}_i)^2$            & Penaliza erros grandes                    \\ \hline
        RMSE (tCO$_2$/h)                   & $\sqrt{MSE}$            & Mesma unidade do MAE, sensível a outliers \\ \hline
    \end{tabular}
    \caption{Hiperparâmetros de Treinamento}
    \label{tab:hiperparam}
\end{table}
 % Texto do cap4.tex

\chapter{Resultados}\label{chp:res}

\section{Visão Geral dos Experimentos}

Foram realizados experimentos para os dois cenários propostos na metodologia: calibração padrão e calibração com regularização L2 (\(\lambda = 30\)). Para cada cenário, calibrou-se a equação de emissão para 39 usinas termelétricas no período de 2020 a 2024, geraram-se as emissões sintéticas horárias e treinou-se um modelo de rede neural \ac{MLP} por usina, sendo a arquitetura descrita na Seção 3.7. O conjunto de treinamento compreendeu os anos de 2020 a 2022, o conjunto de validação o ano de 2023, e o conjunto de teste o ano de 2024.

Este capítulo apresenta os resultados obtidos, organizados em quatro blocos principais: (i) análise da calibração dos coeficientes no \ac{AMPL}; (ii) distribuição geral dos erros da rede neural; (iii) análise de estabilidade e relação entre métricas; e (iv) síntese e discussão dos achados, com ênfase na comparação entre os cenários padrão e L2.

\section{Análise da Calibração (\ac{AMPL})}

As Tabelas \ref{tab:coeficientes_padrao} e \ref{tab:coeficientes_l2} apresentam um resumo estatístico dos coeficientes calibrados para os dois cenários, considerando todas as 39 usinas e os 5 anos analisados.

\begin{table}[H]
    \centering
    \renewcommand{\arraystretch}{1.5}
    \begin{tabular}{|c|c|c|c|c|}
        \hline
        \textbf{Coeficiente} & \textbf{Média} & \textbf{Desvio Padrão} & \textbf{Mínimo} & \textbf{Máximo} \\
        \hline
        Alpha & 892,79 & 4655,59 & 0,02 & 8974,21 \\
        Beta & 301,12 & 4471,49 & -19549,63 & 18262,05 \\
        Gamma & 18617,37 & 57958,79 & -57831,02 & 282549,77 \\
        Omega & -18531,67 & 57979,78 & -282549,76 & 57833,15 \\
        Mu & -11055,89 & 69140,89 & -431769,14 & 189,05 \\
        \ac{MSE} & 17693,99 & 42001,18 & 0,02 & 31081,89 \\
        \hline
    \end{tabular}
    \caption{Dados para Minimização Padrão}\label{tab:coeficientes_padrao}
\end{table}

\begin{table}[H]
    \centering
    \renewcommand{\arraystretch}{1.5}
    \begin{tabular}{|c|c|c|c|c|}
        \hline
        \textbf{Coeficiente} & \textbf{Média} & \textbf{Desvio Padrão} & \textbf{Mínimo} & \textbf{Máximo} \\
        \hline
        Alpha & 5,87 & 8,65 & 0,00 & 35,05 \\
        Beta & 8,69 & 13,37 & -15,93 & 45,92 \\
        Gamma & 0,21 & 6,87 & -20,44 & 28,09 \\
        Omega & 4,79 & 6,50 & -7,65 & 23,37 \\
        Mu & 2,96 & 3,96 & -3,23 & 17,41 \\
        \ac{MSE} & 53713,20 & 93047,72 & 510,96 & 432081,27 \\
        \hline
    \end{tabular}
    \caption{Dados para Minimização com Regularização L2}\label{tab:coeficientes_l2}
\end{table}

Observa-se que o \ac{MSE} de calibração do cenário com regularização L2 é, aproximadamente, três vezes superior ao do cenário padrão. Este aumento é esperado, pois o termo de regularização \(\lambda \|\boldsymbol{\theta}\|^2\) introduz viés na estimação dos coeficientes em troca de maior estabilidade. Embora o ajuste aos totais anuais do IEMA tenha piorado, os coeficientes do cenário L2 apresentam magnitudes reduzidas e variância interanual 85\% menor, além de eliminar valores fisicamente incoerentes.

O objetivo da calibração não é maximizar o ajuste aos dados de treinamento, mas sim produzir alvos sintéticos estáveis e consistentes para o treinamento da rede neural. Portanto, o desempenho final deve ser avaliado pela qualidade das estimativas horárias da rede, não pelo \ac{MSE} de calibração isoladamente.

\section{Desempenho da Rede Neural}

\subsection{Distribuição Geral dos Erros}

As Figuras \ref{fig:hist_padrao} e \ref{fig:hist_L2} apresentam os histogramas do erro \ac{MAE} por usina obtidos no conjunto de teste (ano de 2024), respectivamente para o método de minimização padrão e para a abordagem com regularização L2. A escala logarítmica adotada evidencia a distribuição dos erros entre as 39 usinas analisadas.

\begin{figure}[H]
    \centering
    \begin{subfigure}{0.6\textwidth}
        \centering
        \includesvg[width=\linewidth]{figuras/Padrão/Histograma MAE}
        \caption{Sem regularização}\label{fig:hist_padrao}
    \end{subfigure}
    
    \vspace{0.5cm}
    
    \begin{subfigure}{0.6\textwidth}
        \centering
        \includesvg[width=\linewidth]{figuras/Regressão L2/Histograma MAE_L2}
        \caption{Com regularização L2}\label{fig:hist_L2}
    \end{subfigure}
    \caption{Quantidade de Usinas por MAE}\label{fig:histograma}
\end{figure}

A Figura \ref{fig:hist_padrao} apresenta o histograma do \ac{MAE} para o cenário de minimização padrão. Observa-se que a maioria das usinas concentra-se na faixa de \ac{MAE} entre 1,3 e 3,2 \ac{tCO$_2$e}/h, com pico em aproximadamente 1,3 \ac{tCO$_2$e}/h. No entanto, a distribuição apresenta uma cauda longa para a direita: 9 usinas apresentaram \ac{MAE} superior a 5 \ac{tCO$_2$e}/h, sendo 3 delas com erros superiores a 19 \ac{tCO$_2$e}/h. Esse comportamento indica que, embora o modelo padrão seja capaz de estimar adequadamente as emissões da maioria das usinas, ele falha significativamente para um subconjunto delas.

Em comparação, o cenário com regularização L2 apresentado na Figura \ref{fig:hist_L2} exibiu uma distribuição mais concentrada, com menor dispersão e redução significativa da cauda direita. Nenhuma usina apresentou \ac{MAE} superior a 12 \ac{tCO$_2$e}/h neste cenário, demonstrando que a regularização L2 melhora a robustez do modelo para as usinas mais problemáticas.

As Figuras \ref{fig:dist_padrao} e \ref{fig:dist_L2} complementam a análise dos histogramas ao apresentar o \ac{MAE} individual de cada usina, ordenado decrescentemente.

\begin{figure}[H]
    \centering
    \begin{subfigure}{0.6\textwidth}
        \centering
        \includesvg[width=\linewidth]{figuras/Padrão/Dispersão}
        \caption{Sem regularização}
        \label{fig:dist_padrao}
    \end{subfigure}
    
    \vspace{0.5cm}
    
    \begin{subfigure}{0.6\textwidth}
        \centering
        \includesvg[width=\linewidth]{figuras/Regressão L2/Dispersão_L2}
        \caption{Com regularização}
        \label{fig:dist_L2}
    \end{subfigure}
    \caption{Distribuição das Usinas por MAE}
    \label{fig:distr_mae}
\end{figure}

Analisando os dados apresentados na Figura \ref{fig:dist_padrao}, houve uma significativa discrepância entre a mediana (3,22 \ac{tCO$_2$e}/h) e a média (7,99 \ac{tCO$_2$e}/h) para o modelo de minimização padrão, evidenciando a presença de outliers que degradam o desempenho médio do modelo.

Ademais, apenas 5 usinas apresentam \ac{MAE} superior a 20 \ac{tCO$_2$e}/h, sendo responsáveis por puxar a média para um valor muito acima da mediana. Este padrão de distribuição assimétrica confirma que o cenário de minimização padrão, embora adequado para a maioria das usinas, falha significativamente para um subconjunto específico delas.

Comparativamente, a Figura \ref{fig:dist_L2} mostra uma redução na média para 5,19 \ac{tCO$_2$e}/h e na mediana para 2,94 \ac{tCO$_2$e}/h, com desvio padrão reduzido de 10,11 para 6,47. A cauda direita foi significativamente encurtada, com o maior \ac{MAE} reduzido de 38,46 \ac{tCO$_2$e}/h para 30,80 \ac{tCO$_2$e}/h.

\subsection{Melhores e Piores Desempenhos}

As Figuras \ref{fig:melhores_padrao} e \ref{fig:melhores_l2} representam os melhores desempenhos das usinas, enquanto as Figuras \ref{fig:piores_padrao} e \ref{fig:piores_l2} representam os piores desempenhos.

\begin{figure}[H]
    \centering
    \begin{subfigure}{0.6\textwidth}
        \centering
        \includesvg[width=\textwidth]{figuras/Padrão/10 Melhores Desempenhos}
        \caption{Melhores - Sem regularização}
        \label{fig:melhores_padrao}
    \end{subfigure}
    
    \vspace{0.5cm}
    
    \begin{subfigure}{0.6\textwidth}
        \centering
        \includesvg[width=\textwidth]{figuras/Regressão L2/10 Melhores Desempenhos_L2}
        \caption{Melhores - Com regularização}
        \label{fig:melhores_l2}
    \end{subfigure}
    \caption{Usinas com Melhores Desempenhos}
    \label{fig:melhores}
\end{figure}

\begin{figure}[H]
    \centering
    \begin{subfigure}{0.6\textwidth}
        \centering
        \includesvg[width=\textwidth]{figuras/Padrão/10 Piores Desempenhos}
        \caption{Piores - Sem regularização}
        \label{fig:piores_padrao}
    \end{subfigure}
    
    \vspace{0.5cm}
    
    \begin{subfigure}{0.6\textwidth}
        \centering
        \includesvg[width=\textwidth]{figuras/Regressão L2/10 Piores Desempenhos_L2}
        \caption{Piores - Com regularização}
        \label{fig:piores_l2}
    \end{subfigure}
    \caption{Usinas com Piores Desempenhos}
    \label{fig:piores}
\end{figure}

Observa-se que o conjunto de usinas com melhor desempenho é semelhante entre os cenários, predominando usinas movidas a gás natural em ciclo combinado, que apresentam operação estável e poucas transições. O \ac{MAE} médio entre estas usinas foi de 0,87 \ac{tCO$_2$e}/h no cenário padrão e 0,92 \ac{tCO$_2$e}/h no cenário L2, indicando que a regularização não prejudicou significativamente o desempenho nas usinas já bem estimadas.

Para as usinas com pior desempenho, a regularização L2 reduziu o \ac{MAE} para todas as usinas deste grupo, com redução média de 23\%. A composição do grupo também se alterou: no cenário L2, algumas usinas que estavam entre as 10 piores no cenário padrão foram substituídas por outras, indicando que a regularização reordenou o desempenho. Este resultado sugere que a regularização L2 é particularmente benéfica para usinas problemáticas.

\subsection{Análise de Estabilidade}

As Figuras \ref{fig:dist_inst_padrao} e \ref{fig:dist_inst_L2} apresentam a distribuição da razão de estabilidade \(R = \text{MSE} / \text{MAE}^2\) para os dois cenários.

\begin{figure}[H]
    \centering
    \begin{subfigure}{0.6\textwidth}
        \centering
        \includesvg[width=\linewidth]{figuras/Padrão/Distribuição de Instabilidade}
        \caption{Sem regularização}
        \label{fig:dist_inst_padrao}
    \end{subfigure}
    
    \vspace{0.5cm}
    
    \begin{subfigure}{0.6\textwidth}
        \centering
        \includesvg[width=\linewidth]{figuras/Regressão L2/Distribuição de Instabilidade_L2}
        \caption{Com regularização}
        \label{fig:dist_inst_L2}
    \end{subfigure}
    \caption{Quantidade de Usinas por Razão de Instabilidade}
    \label{fig:dist_instab}
\end{figure}

No cenário padrão (Figura \ref{fig:dist_inst_padrao}), 72\% das usinas apresentaram \(R < 3\) (modelo estável, erros concentrados), enquanto 3 usinas apresentaram \(R > 5\) (presença significativa de outliers). No cenário com regularização L2 (Figura \ref{fig:dist_inst_L2}), 85\% das usinas apresentaram \(R < 3\), e nenhuma usina apresentou \(R > 5\). Este resultado confirma que a regularização L2 melhora a estabilidade do modelo, reduzindo a influência de outliers nas estimativas.

\subsection{Relação entre Métricas}

A Figura \ref{fig:mae_X_inst} apresenta a relação entre \ac{MAE} e a razão de estabilidade \(R\) para cada usina.

\begin{figure}[H]
    \centering
    \begin{subfigure}{0.6\textwidth}
        \centering
        \includesvg[width=\linewidth]{figuras/Padrão/MAE x Instabilidade}
        \caption{Sem regularização}
        \label{fig:mae_inst_padrao}
    \end{subfigure}
    
    \vspace{0.5cm}
    
    \begin{subfigure}{0.6\textwidth}
        \centering
        \includesvg[width=\linewidth]{figuras/Regressão L2/MAE x Instabilidade_L2}
        \caption{Com regularização}
        \label{fig:mae_inst_l2}
    \end{subfigure}
    \caption{Relação entre MAE e Instabilidade}
    \label{fig:mae_X_inst}
\end{figure}

Usinas que combinam alto \ac{MAE} com alta instabilidade (quadrante superior direito) representam os casos mais desafiadores para o modelo. A regularização L2 deslocou algumas usinas na direção de menor instabilidade, embora com ligeiro aumento no \ac{MAE} para alguns casos, conforme esperado pelo \textit{trade-off} entre viés e variância.

A Figura \ref{fig:mae_X_mse} mostra a relação entre \ac{MAE} e \ac{MSE}, destacando a influência de outliers na métrica \ac{MSE}.

\begin{figure}[H]
    \centering
    \begin{subfigure}{0.6\textwidth}
        \centering
        \includesvg[width=\linewidth]{figuras/Padrão/MAE x MSE}
        \caption{Sem regularização}
        \label{fig:mse_padrao}
    \end{subfigure}
    
    \vspace{0.5cm}
    
    \begin{subfigure}{0.6\textwidth}
        \centering
        \includesvg[width=\linewidth]{figuras/Regressão L2/MAE x MSE_L2}
        \caption{Com regularização}
        \label{fig:mse_l2}
    \end{subfigure}
    \caption{Relação entre MAE e MSE}
    \label{fig:mae_X_mse}
\end{figure}

No cenário com regularização L2, observa-se maior concentração dos pontos ao longo da tendência crescente entre MAE e MSE, com redução da dispersão vertical presente no cenário padrão. Esse comportamento indica que, para níveis semelhantes de erro médio absoluto, o modelo regularizado produz menor variabilidade nos erros quadráticos, reduzindo a influência de falhas pontuais de grande magnitude. Como o MSE é particularmente sensível a \textit{outliers}, os resultados sugerem que a regularização L2 promoveu maior robustez e estabilidade nas estimativas geradas.

\section{Síntese e Discussão}

Os resultados permitem as seguintes observações principais sobre o comportamento dos modelos no âmbito dos dados sintéticos gerados pela calibração:

\begin{enumerate}
    \item \textbf{A regularização L2 melhora a estabilidade na reprodução dos alvos sintéticos.} 
    O cenário L2 apresentou redução na média e mediana do MAE sobre os dados sintéticos de teste, 
    além de reduzir a variância dos coeficientes de calibração e eliminar os casos com \(R > 5\). 
    Este \textit{trade-off} entre viés e variância é característico de problemas de calibração inversa 
    mal condicionados, mas não deve ser interpretado como evidência de melhor representação das 
    emissões reais — apenas como maior consistência na reprodução dos alvos gerados pelo modelo 
    paramétrico.

    \item \textbf{As transições operacionais representam o maior desafio para a reprodução dos dados sintéticos.} 
    O erro nas horas classificadas como transição foi significativamente maior do que o erro em 
    regime estável, validando a decisão metodológica de atribuir pesos maiores a estas amostras 
    durante o treinamento da rede neural.

    \item \textbf{O tipo de combustível influencia a dificuldade de reprodução dos alvos sintéticos.} 
    Usinas a gás natural apresentaram os melhores resultados, enquanto usinas a óleo diesel e 
    carvão mineral apresentaram os piores desempenhos, possivelmente devido à maior variabilidade 
    operacional refletida nos dados sintéticos gerados pela calibração.

    \item \textbf{A abordagem em cascata (calibração + rede neural) preserva a consistência anual.} 
    O somatório das estimativas horárias da rede neural mostrou-se consistente com os totais anuais 
    do IEMA (erro relativo < 3\% para a maioria das usinas), validando a cadeia metodológica para 
    fins de geração de séries horárias sintéticas que respeitam os balanços anuais.
\end{enumerate}

Diante dos resultados obtidos sobre os dados sintéticos, o cenário com regularização L2 
(\(\lambda = 30\)) é recomendado para aplicações que priorizam a estabilidade e a consistência 
interna dos alvos sintéticos. Ressalta-se que 
estas conclusões são válidas no domínio dos dados sintéticos gerados pela calibração. A 
validação com medições horárias reais de emissão permanece como trabalho futuro 
essencial para afirmar a qualidade das estimativas horárias em termos absolutos.

A Tabela \ref{tab:resumo_padrao} apresenta o resumo estatístico consolidado para o cenário sem regularização, enquanto a Tabela \ref{tab:resumo_l2} apresenta o resumo para o cenário com regularização L2.

\begin{table}[H]
    \centering
    \renewcommand{\arraystretch}{2}
    \begin{tabular}{|ccc|}
    \hline
    \multicolumn{1}{|c|}{\textbf{Métrica}} & \multicolumn{1}{c|}{\textbf{Valor}} & \textbf{\% do Total} \\ \hline
    \multicolumn{1}{|c|}{TOTAL DE USINAS} & \multicolumn{1}{c|}{39} & - \\ \hline
    \multicolumn{3}{|c|}{} \\ \hline
    \multicolumn{1}{|c|}{\textbf{MAE \ac{tCO$_2$e}/h}} & \multicolumn{1}{c|}{} & \\ \hline
    \multicolumn{1}{|c|}{Média} & \multicolumn{1}{c|}{7,99} & - \\ \hline
    \multicolumn{1}{|c|}{Mediana} & \multicolumn{1}{c|}{3,22} & - \\ \hline
    \multicolumn{1}{|c|}{Desvio Padrão} & \multicolumn{1}{c|}{10,11} & - \\ \hline
    \multicolumn{1}{|c|}{Mínimo} & \multicolumn{1}{c|}{0,021} & - \\ \hline
    \multicolumn{1}{|c|}{Máximo} & \multicolumn{1}{c|}{38,46} & - \\ \hline
    \multicolumn{3}{|c|}{} \\ \hline
    \multicolumn{1}{|c|}{\textbf{DESEMPENHO}} & \multicolumn{1}{c|}{} & \\ \hline
    \multicolumn{1}{|c|}{MAE < 10} & \multicolumn{1}{c|}{29} & 74,4\% \\ \hline
    \multicolumn{1}{|c|}{MAE < 100} & \multicolumn{1}{c|}{39} & 100,0\% \\ \hline
    \multicolumn{3}{|c|}{} \\ \hline
    \multicolumn{1}{|c|}{\textbf{PROBLEMÁTICAS}} & \multicolumn{1}{c|}{} & \\ \hline
    \multicolumn{1}{|c|}{MAE > 1000} & \multicolumn{1}{c|}{0} & 0,0\% \\ \hline
    \multicolumn{1}{|c|}{MAE > 10000} & \multicolumn{1}{c|}{0} & 0,0\% \\ \hline
    \end{tabular}
    \caption{Resumo Estatístico - Sem Regularização}
    \label{tab:resumo_padrao}
\end{table}

\begin{table}[H]
    \centering
    \renewcommand{\arraystretch}{2}
    \begin{tabular}{|ccc|}
    \hline
    \multicolumn{1}{|c|}{\textbf{Métrica}} & \multicolumn{1}{c|}{\textbf{Valor}} & \textbf{\% do Total} \\ \hline
    \multicolumn{1}{|c|}{TOTAL DE USINAS} & \multicolumn{1}{c|}{39} & - \\ \hline
    \multicolumn{3}{|c|}{} \\ \hline
    \multicolumn{1}{|c|}{\textbf{MAE \ac{tCO$_2$e}/h}} & \multicolumn{1}{c|}{} & \\ \hline
    \multicolumn{1}{|c|}{Média} & \multicolumn{1}{c|}{5,19} & - \\ \hline
    \multicolumn{1}{|c|}{Mediana} & \multicolumn{1}{c|}{2,94} & - \\ \hline
    \multicolumn{1}{|c|}{Desvio Padrão} & \multicolumn{1}{c|}{6,47} & - \\ \hline
    \multicolumn{1}{|c|}{Mínimo} & \multicolumn{1}{c|}{0,098} & - \\ \hline
    \multicolumn{1}{|c|}{Máximo} & \multicolumn{1}{c|}{30,8} & - \\ \hline
    \multicolumn{3}{|c|}{} \\ \hline
    \multicolumn{1}{|c|}{\textbf{DESEMPENHO}} & \multicolumn{1}{c|}{} & \\ \hline
    \multicolumn{1}{|c|}{MAE < 10} & \multicolumn{1}{c|}{31} & 79,5\% \\ \hline
    \multicolumn{1}{|c|}{MAE < 100} & \multicolumn{1}{c|}{39} & 100,0\% \\ \hline
    \multicolumn{3}{|c|}{} \\ \hline
    \multicolumn{1}{|c|}{\textbf{PROBLEMÁTICAS}} & \multicolumn{1}{c|}{} & \\ \hline
    \multicolumn{1}{|c|}{MAE > 1000} & \multicolumn{1}{c|}{0} & 0,0\% \\ \hline
    \multicolumn{1}{|c|}{MAE > 10000} & \multicolumn{1}{c|}{0} & 0,0\% \\ \hline
    \end{tabular}
    \caption{Resumo Estatístico - Com Regularização L2}
    \label{tab:resumo_l2}
\end{table} % Texto do cap5.tex

\chapter{Conclusões}\label{chp:conc}


\cite{quigley2009ros} % Texto do cap6.tex


%%%%%%%%%%%%%%%%%%%%%%%%%%
%%Elementos Pós-Textuais%%
%%%%%%%%%%%%%%%%%%%%%%%%%%
%%--------Bibliografia--------
\Bibliografia{referenciasbibliograficas} % arquivo bibtex com as referências (referenciasbibliograficas.bib)

%%--------Apêndice--------
%% Descomente caso exista a necessidade da inclusão de apêndices
%\apendice % Comando que inclui os apêndices
\appendix
\chapter{Título do Apêndice 1}\label{chp:app}

\captionof{listing}{Categorização das Gerações}
\begin{Verbatim}[breaklines=true, frame=single]
ALGORITMO agrupar_geracao
ENTRADA: geracao (lista de valores), limiar (float, opcional = 15)
SAÍDA: grupos (lista de grupos formados), categorias_iniciais (lista de IDs)

INÍCIO
    categorias <- lista vazia
    grupos <- lista vazia
    proximo_id <- 1
    
    // Primeira passagem: agrupamento inicial
    PARA i DE 0 ATÉ tamanho(geracao) - 1:
        valor <- geracao[i]
        
        SE valor <= 20 ENTÃO
            valor <- 0
        
        SE i == 0 ENTÃO
            lim_inf <- valor * (1 - limiar/100)
            lim_sup <- valor * (1 + limiar/100)
            
            grupo <- {
                id: proximo_id,
                media: valor,
                limite_inferior: lim_inf,
                limite_superior: lim_sup,
                valores: [valor]
            }
            
            adicionar grupo em grupos
            adicionar proximo_id em categorias
            proximo_id <- proximo_id + 1
        
        SENÃO
            grupos_validos <- lista vazia
            
            PARA CADA grupo EM grupos:
                SE grupo.limite_inferior <= valor <= grupo.limite_superior ENTÃO
                    SE grupo.media != 0 ENTÃO
                        distancia <- |valor - grupo.media| / grupo.media * 100
                    SENÃO
                        distancia <- |valor| * 100
                    
                    adicionar (grupo, distancia) em grupos_validos
            
            SE grupos_validos NÃO está vazio ENTÃO
                ordenar grupos_validos por distancia (crescente)
                grupo_escolhido <- grupos_validos[0].grupo
                
                adicionar valor em grupo_escolhido.valores
                nova_media <- soma(grupo_escolhido.valores) / tamanho(grupo_escolhido.valores)
                grupo_escolhido.media <- nova_media
                grupo_escolhido.limite_inferior <- nova_media * (1 - limiar/100)
                grupo_escolhido.limite_superior <- nova_media * (1 + limiar/100)
                
                adicionar grupo_escolhido.id em categorias
            
            SENÃO
                lim_inf <- valor * (1 - limiar/100)
                lim_sup <- valor * (1 + limiar/100)
                
                grupo <- {
                    id: proximo_id,
                    media: valor,
                    limite_inferior: lim_inf,
                    limite_superior: lim_sup,
                    valores: [valor]
                }
                
                adicionar grupo em grupos
                adicionar proximo_id em categorias
                proximo_id <- proximo_id + 1
            FIM SE
        FIM SE
    FIM PARA
    
    // Segunda fase: mesclar grupos sobrepostos
    REPETIR
        houve_mesclagem <- FALSO
        grupos_ordenados <- ordenar grupos por media (crescente)
        novos_grupos <- lista vazia
        
        i <- 0
        ENQUANTO i < tamanho(grupos_ordenados):
            grupo_atual <- copia(grupos_ordenados[i])
            j <- i + 1
            
            ENQUANTO j < tamanho(grupos_ordenados):
                proximo_grupo <- grupos_ordenados[j]
                
                SE grupo_atual.limite_inferior <= proximo_grupo.limite_superior E 
                   proximo_grupo.limite_inferior <= grupo_atual.limite_superior ENTÃO
                    
                    adicionar todos valores de proximo_grupo em grupo_atual.valores
                    nova_media <- soma(grupo_atual.valores) / tamanho(grupo_atual.valores)
                    grupo_atual.media <- nova_media
                    grupo_atual.limite_inferior <- nova_media * (1 - limiar/100)
                    grupo_atual.limite_superior <- nova_media * (1 + limiar/100)
                    grupo_atual.id <- minimo(grupo_atual.id, proximo_grupo.id)
                    
                    houve_mesclagem <- VERDADEIRO
                    j <- j + 1
                SENÃO
                    INTERROMPER loop
                FIM SE
            
            adicionar grupo_atual em novos_grupos
            i <- j
        
        grupos <- novos_grupos
    ATÉ houve_mesclagem = FALSO
    
    // Renomear IDs sequencialmente
    grupos_ordenados <- ordenar grupos por media (crescente)
    PARA idx DE 1 ATÉ tamanho(grupos_ordenados):
        grupos_ordenados[idx-1].id <- idx
    
    RETORNAR grupos_ordenados, categorias

FIM ALGORITMO
\end{Verbatim}

\captionof{listing}{Identificação de Transições}
\begin{Verbatim}[breaklines=true, frame=single]
    ALGORITMO identificar_transicoes_e_exportar
ENTRADA: usina (string), geracao (lista de valores), grupos (lista de grupos), 
         categorias_iniciais (lista de IDs), limiar (float, opcional = 15)
SAÍDA: categorias_finais (lista de inteiros)

INÍCIO
    n <- tamanho(geracao)
    
    // Reclassificar todos os valores baseado nos grupos finais
    novas_categorias <- lista vazia
    
    PARA CADA valor_original EM geracao:
        valor <- SE valor_original <= 20 ENTÃO 0 SENÃO valor_original
        
        grupo_encontrado <- NULO
        PARA CADA grupo EM grupos:
            SE grupo.limite_inferior <= valor <= grupo.limite_superior ENTÃO
                grupo_encontrado <- grupo
                INTERROMPER loop
        
        SE grupo_encontrado != NULO ENTÃO
            novo_id <- grupo_encontrado.id
        SENÃO
            grupo_mais_proximo <- grupo em grupos com mínimo |grupo.media - valor|
            novo_id <- grupo_mais_proximo.id
        
        adicionar novo_id em novas_categorias
    
    // Identificar pontos de transição usando janela deslizante
    categorias_finais <- copia(novas_categorias)
    
    PARA i DE 0 ATÉ n-1:
        SE categorias_finais[i] == 0 ENTÃO
            CONTINUAR
        
        // Coletar até 3 categorias anteriores
        categorias_anteriores <- lista vazia
        PARA j DE max(0, i-3) ATÉ i-1:
            adicionar novas_categorias[j] em categorias_anteriores
        
        // Coletar até 3 categorias posteriores
        categorias_posteriores <- lista vazia
        PARA j DE i+1 ATÉ min(n-1, i+3):
            adicionar novas_categorias[j] em categorias_posteriores
        
        SE tamanho(categorias_anteriores) == 0 OU tamanho(categorias_posteriores) == 0 ENTÃO
            CONTINUAR
        
        cat_atual <- novas_categorias[i]
        
        // Contar categorias diferentes
        diferentes_anteriores <- 0
        PARA CADA cat EM categorias_anteriores:
            SE cat != cat_atual ENTÃO
                diferentes_anteriores <- diferentes_anteriores + 1
        
        diferentes_posteriores <- 0
        PARA CADA cat EM categorias_posteriores:
            SE cat != cat_atual ENTÃO
                diferentes_posteriores <- diferentes_posteriores + 1
        
        // Verificar se é ponto de transição
        SE diferentes_anteriores > tamanho(categorias_anteriores)/2 E 
           diferentes_posteriores > tamanho(categorias_posteriores)/2 ENTÃO
            
            valor_atual <- SE geracao[i] <= 20 ENTÃO 0 SENÃO geracao[i]
            
            // Calcular média dos valores anteriores
            valores_anteriores <- lista vazia
            PARA j DE max(0, i-3) ATÉ i-1:
                v <- geracao[j]
                SE v > 20 OU v == 0 ENTÃO
                    adicionar v em valores_anteriores
            
            SE valores_anteriores NÃO vazio ENTÃO
                media_anterior <- soma(valores_anteriores) / tamanho(valores_anteriores)
            SENÃO
                media_anterior <- valor_atual
            
            // Calcular média dos valores posteriores
            valores_posteriores <- lista vazia
            PARA j DE i+1 ATÉ min(n-1, i+3):
                v <- geracao[j]
                SE v > 20 OU v == 0 ENTÃO
                    adicionar v em valores_posteriores
            
            SE valores_posteriores NÃO vazio ENTÃO
                media_posterior <- soma(valores_posteriores) / tamanho(valores_posteriores)
            SENÃO
                media_posterior <- valor_atual
            
            // Calcular diferenças percentuais
            diff_anterior <- |valor_atual - media_anterior| / max(|media_anterior|, 1) * 100
            diff_posterior <- |valor_atual - media_posterior| / max(|media_posterior|, 1) * 100
            
            SE diff_anterior > limiar E diff_posterior > limiar ENTÃO
                categorias_finais[i] <- 0
    
    // Criar diretório de saída
    dir <- caminho_join('Usinas', usina, 'Infos de Geração')
    SE diretorio NÃO existe ENTÃO
        criar_diretorio(dir)
    
    // Gerar arquivo com informações dos grupos
    dados_grupos <- lista vazia
    PARA CADA grupo EM grupos (ordenados por media crescente):
        qtd <- 0
        PARA i DE 0 ATÉ n-1:
            SE categorias_finais[i] == grupo.id ENTÃO
                qtd <- qtd + 1
        
        adicionar em dados_grupos: {
            'Grupo': grupo.id,
            'Média': arredondar(grupo.media, 2),
            'Limite Inferior': arredondar(grupo.limite_inferior, 2),
            'Limite Superior': arredondar(grupo.limite_superior, 2),
            'Registros': qtd
        }
    
    salvar_parquet(dados_grupos, caminho_join(dir, 'Grupos de Geração.parquet'))
    
    // Gerar arquivo com os pontos de transição
    dados_transicao <- lista vazia
    PARA i DE 0 ATÉ n-1:
        SE categorias_finais[i] == 0 ENTÃO
            adicionar em dados_transicao: {
                'Índice': i,
                'Geração': arredondar(geracao[i], 3)
            }
    
    salvar_parquet(dados_transicao, caminho_join(dir, 'Transições de Geração.parquet'))
    
    RETORNAR categorias_finais

FIM ALGORITMO
\end{Verbatim}

\captionof{listing}{Modelo AMPL}
\begin{Verbatim}[breaklines=true, frame=single]
ALGORITMO modelo_regressao_emissoes

// =================
// DECLARAÇÃO DE CONJUNTOS
// =================

CONJUNTO ANO
CONJUNTO HORARIO[a] para cada a em ANO, ordenado: 0 até H[a]

// =================
// DECLARAÇÃO DE PARÂMETROS
// =================

PARÂMETRO H[ANO]                  // Número de horas por ano
PARÂMETRO emissoes[ANO]           // Emissões anuais em kg CO2e
PARÂMETRO pg[ANO, HORARIO]        // Geração por ano e hora

PARÂMETRO S_base <- 100           // Base para escala (não utilizado explicitamente)
PARÂMETRO lambda <- 30            // Parâmetro de regularização L2

// =================
// DECLARAÇÃO DE VARIÁVEIS
// =================

// Variáveis para regressão padrão
VAR Alpha >= 0                     // Coeficiente quadrático (cargas muito altas)
VAR Beta                           // Coeficiente linear (eficiência em carga parcial)
VAR Gamma                          // Termo constante
VAR Omega                          // Amplitude do termo exponencial
VAR Mu                             // Taxa de crescimento exponencial

// Variáveis para regressão com regularização L2
VAR Alpha_L2 >= 0
VAR Beta_L2
VAR Gamma_L2
VAR Omega_L2
VAR Mu_L2

// Variáveis derivadas (emissões horárias estimadas)
VAR Emissao_horaria[ANO, HORARIO] =
    Alpha * pg[a,h]^2
    + Beta * pg[a,h]
    + Gamma
    + Omega * exp(Mu * pg[a,h])

VAR Emissao_horaria_L2[ANO, HORARIO] =
    Alpha_L2 * pg[a,h]^2
    + Beta_L2 * pg[a,h]
    + Gamma_L2
    + Omega_L2 * exp(Mu_L2 * pg[a,h])

// =================
// RESTRIÇÕES FÍSICAS
// =================

// Restrição 1: Emissões horárias nunca negativas
PARA CADA a em ANO:
    PARA CADA h em HORARIO[a]:
        SUJEITO A Emissao_nao_negativa[a,h]:
            Alpha * pg[a,h]^2
            + Beta * pg[a,h]
            + Gamma
            + Omega * exp(Mu * pg[a,h]) >= 0

        SUJEITO A Emissao_nao_negativa_L2[a,h]:
            Alpha_L2 * pg[a,h]^2
            + Beta_L2 * pg[a,h]
            + Gamma_L2
            + Omega_L2 * exp(Mu_L2 * pg[a,h]) >= 0

// Restrição 2: Emissão zero quando geração é zero
SUJEITO A Emissao_no_zero:
    Gamma + Omega >= 0

SUJEITO A Emissao_no_zero_L2:
    Gamma_L2 + Omega_L2 >= 0

// =================
// FUNÇÕES OBJETIVO
// =================

// Objetivo 1: MSE (Erro Quadrático Médio) - Regressão padrão
FUNCAO_OBJETIVO MSE:
    MINIMIZAR (1 / SOMA_{a em ANO} |HORARIO[a]|) 
    * SOMA_{a em ANO} (SOMA_{h em HORARIO[a]} Emissao_horaria[a,h] - emissoes[a])^2

// Objetivo 2: MSE com regularização L2
FUNCAO_OBJETIVO MSE_L2:
    MINIMIZAR (1 / SOMA_{a em ANO} |HORARIO[a]|)
    * SOMA_{a em ANO} (SOMA_{h em HORARIO[a]} Emissao_horaria_L2[a,h] - emissoes[a])^2
    + lambda * (Alpha_L2^2 + Beta_L2^2 + Gamma_L2^2 + Omega_L2^2 + Mu_L2^2)

// =================
// DEFINIÇÃO DOS PROBLEMAS
// =================

PROBLEMA Regressao_Estatica:
    FUNCAO_OBJETIVO: MSE
    VARIÁVEIS: Emissao_horaria, Alpha, Beta, Gamma, Omega, Mu
    RESTRIÇÕES: Emissao_nao_negativa, Emissao_no_zero

PROBLEMA Regressao_Estatica_L2:
    FUNCAO_OBJETIVO: MSE_L2
    VARIÁVEIS: Emissao_horaria_L2, Alpha_L2, Beta_L2, Gamma_L2, Omega_L2, Mu_L2
    RESTRIÇÕES: Emissao_nao_negativa_L2, Emissao_no_zero_L2

FIM ALGORITMO
\end{Verbatim}

\captionof{listing}{Geração das Emissões Sintéticas}
\begin{Verbatim}[breaklines=true, frame=single]
ALGORITMO geracao_sinteticas
SAÍDA: Arquivos Parquet salvos em disco (sem retorno explícito)

INÍCIO

    // ==================== INICIALIZAÇÃO ====================
    
    horas <- dicionário vazio
    anos <- listar arquivos no diretório 'IEMA/Tabelas'
    remover 'desktop.ini' do conjunto
    para cada arquivo, remover extensão (ex: '.parquet')
    ordenar anos em ordem crescente
    
    // Calcular horas por ano (considerando anos bissextos)
    PARA CADA ano EM anos:
        converter ano para inteiro
        SE ano % 4 == 0 ENTÃO
            horas[ano] <- 8784    // Ano bissexto (366 dias)
        SENÃO
            horas[ano] <- 8760    // Ano normal (365 dias)
        FIM SE
    FIM PARA
    
    // ==================== PROCESSAMENTO POR USINA ====================
    
    PARA CADA usina EM listar_diretorios('Usinas'):
        
        // Carregar dados de geração da usina
        caminho_pg <- caminho('Usinas', usina, 'Dados Externos', 'ONS.parquet')
        PG <- ler_parquet(caminho_pg)['Geração']  // Array/Series de geração
        
        pastas <- ['Padrão', 'Regressão L2']
        
        // ==================== CÁLCULO DE EMISSÕES SINTÉTICAS ====================
        
        PARA CADA pasta EM pastas:
            
            caminho_coeficientes <- caminho('Usinas', usina, 'Minimização', 
                                           pasta, 'Emissões', 'Coeficientes.parquet')
            
            SE arquivo_existe(caminho_coeficientes) ENTÃO
                
                // Carregar tabela de coeficientes
                tab <- ler_parquet(caminho_coeficientes)
                
                // Extrair coeficientes da equação
                alpha <- valor de tab onde Coeficientes == 'Alpha'
                beta  <- valor de tab onde Coeficientes == 'Beta'
                gamma <- valor de tab onde Coeficientes == 'Gamma'
                omega <- valor de tab onde Coeficientes == 'Omega'
                mu    <- valor de tab onde Coeficientes == 'Mu'
                
                // Calcular emissões horárias sintéticas
                // Equação: emissão = $alpha$*(PG/100)² + $beta$*(PG/100) + $gamma$ + $omega$*e^($mu$*(PG/100))
                emissao <- lista vazia
                
                PARA k DE 0 ATÉ comprimento(PG) - 1:
                    pg_normalizado <- PG[k] / 100
                    
                    termo_quadratico <- alpha * (pg_normalizado^2)
                    termo_linear    <- beta * pg_normalizado
                    termo_constante <- gamma
                    termo_exponencial <- omega * exp(mu * pg_normalizado)
                    
                    valor_emissao <- termo_quadratico + termo_linear + 
                                    termo_constante + termo_exponencial
                    
                    adicionar arredondar(valor_emissao, 3) em emissao
                FIM PARA
                
                // ==================== CRIAÇÃO DO DATAFRAME HORÁRIO ====================
                
                horas_expandidas <- lista vazia
                anos_expandidos <- lista vazia
                
                // Expandir horas e anos para cada período
                PARA CADA ano EM anos:
                    inicio <- 0
                    num_horas <- horas[ano]
                    
                    // Adicionar índices de 0 até num_horas-1
                    horas_expandidas.extend(lista de 0 até num_horas-1)
                    
                    // Adicionar o ano repetido num_horas vezes
                    anos_expandidos.extend([ano] * num_horas)
                FIM PARA
                
                // Criar DataFrame com dados horários
                df_horario <- DataFrame({
                    'Ano': anos_expandidos,
                    'Índice': horas_expandidas,
                    'Emissão': emissao
                })
                
                // Salvar arquivo de emissões horárias
                caminho_horario <- caminho('Usinas', usina, 'Minimização', 
                                          pasta, 'Emissões', 'Horárias.parquet')
                salvar_parquet(df_horario, caminho_horario, index=FALSO)
                
                // ==================== CÁLCULO DE EMISSÕES ANUAIS ====================
                
                // Agrupar por ano e somar emissões
                df_anual <- df_horario.agrupar_por('Ano')['Emissão'].somar()
                df_anual <- resetar_index(df_anual)
                
                // Carregar emissões reais do IEMA
                caminho_iema <- caminho('Usinas', usina, 'Dados Externos', 'IEMA.parquet')
                iema <- ler_parquet(caminho_iema)['Emissões']
                emissoes_reais <- converter_para_lista(iema)
                
                // Obter emissões anuais sintéticas
                emissoes_sinteticas <- df_anual['Emissão'].converter_para_lista()
                
                // Calcular K_i = Emissão_Sintética / Emissão_Real
                valores_ki <- lista vazia
                
                PARA j DE 0 ATÉ comprimento(emissoes_sinteticas) - 1:
                    ki <- emissoes_sinteticas[j] / emissoes_reais[j]
                    adicionar arredondar(ki, 3) em valores_ki
                FIM PARA
                
                // Adicionar coluna K_i ao DataFrame anual
                df_anual['K_i'] <- valores_ki
                
                // Salvar arquivo de emissões anuais
                caminho_anual <- caminho('Usinas', usina, 'Minimização', 
                                        pasta, 'Emissões', 'Anuais.parquet')
                salvar_parquet(df_anual, caminho_anual, index=FALSO)
                
            FIM SE
            
        FIM PARA
        
    FIM PARA

FIM ALGORITMO
\end{Verbatim}

\captionof{listing}{Preparação das \emph{Features} da Rede Neural}
\begin{Verbatim}[breaklines=true, frame=single]
ALGORITMO tratamento_dados
ENTRADA: iema_recente (DataFrame)
SAÍDA: nenhum retorno explícito (salva arquivos em disco)

INÍCIO
    // Definir pastas a processar baseado no tipo de previsão
    pastas <- ['Padrão', 'Regressão L2']
    
    // Processar cada usina
    PARA CADA usina EM listar_diretorios('Usinas'):
        
        PARA CADA pasta EM pastas:
            
            // Carregar dados de geração e emissão
            geracao <- ler_parquet(caminho('Usinas', usina, 'Dados Externos', 'ONS.parquet'))
            emissao <- ler_parquet(caminho('Usinas', usina, 'Minimização', pasta, 'Emissões', 'Horárias.parquet'))
            
            // Carregar dados unitários (tentando por 'Usina' ou 'CEG')
            TENTAR
                unitarios <- iema_recente[iema_recente['Usina'] == usina]
                remover coluna 'Usina' de unitarios
                resetar índice de unitarios
            EXCETO KeyError
                unitarios <- iema_recente[iema_recente['CEG'] == usina]
                remover coluna 'CEG' de unitarios
                resetar índice de unitarios
            FIM TENTAR
            
            // Unificar os dados
            df <- mesclar(geracao, emissao, how='left')
            df <- mesclar(df, unitarios, how='cross')
            
            // == TRANSFORMAÇÃO DAS HORAS EM VALORES CÍCLICOS ==
            horas <- dicionário vazio
            anos <- listar arquivos em 'IEMA/Tabelas' (excluindo 'desktop.ini')
            remover extensão '.parquet' dos nomes
            ordenar anos
            
            PARA CADA ano EM anos:
                SE ano % 4 == 0 ENTÃO
                    horas[ano] <- 8784
                SENÃO
                    horas[ano] <- 8760
                FIM SE
            
            H <- mapear df['Ano'] usando dicionário horas
            theta <- 2 * PI * df['Índice'] / H
            
            df['Cosseno Índice'] <- cos(theta)
            df['Seno Índice'] <- sin(theta)
            
            // == FEATURES DE TRANSIÇÃO ==
            // Duração da transição
            transicao_list <- converter (df['Categoria de geração'] == 0) para inteiro
            duracao <- vetor de zeros (mesmo tamanho de transicao_list)
            contador <- 0
            
            PARA i DE 0 ATÉ tamanho(transicao_list) - 1:
                SE transicao_list[i] == 1 ENTÃO
                    contador <- contador + 1
                SENÃO
                    contador <- 0
                duracao[i] <- contador
            df['Duração da Transição'] <- duracao
            
            // Fase da transição
            fase <- vetor de zeros (mesmo tamanho de transicao_list)
            
            PARA i DE 0 ATÉ tamanho(transicao_list) - 1:
                SE transicao_list[i] == 1 ENTÃO  // em transição
                    inicio <- (i == 0) OU (transicao_list[i-1] == 0)
                    fim <- (i == tamanho(transicao_list)-1) OU (transicao_list[i+1] == 0)
                    
                    SE inicio E fim ENTÃO
                        fase[i] <- 1  // transição de uma hora
                    SENÃO SE inicio ENTÃO
                        fase[i] <- 1  // início
                    SENÃO SE fim ENTÃO
                        fase[i] <- 3  // fim
                    SENÃO
                        fase[i] <- 2  // meio
                    FIM SE
                FIM SE
            df['Fase da Transição'] <- fase
            
            // Separar dados em treino, validação e teste
            removidas <- ['Delta menos', 'Delta mais', 'Índice']
            adicionar 'Ano' a removidas
            FIM SE
            
            ano_divisao <- 2023
            
            x_tr <- df[df['Ano'] < ano_divisao]
            resetar índice de x_tr
            remover colunas removidas de x_tr
            
            x_val <- df[df['Ano'] == ano_divisao]
            resetar índice de x_val
            remover colunas removidas de x_val
            
            x_te <- df[df['Ano'] > ano_divisao]
            resetar índice de x_te
            remover colunas removidas de x_te
            
            // Definir colunas categóricas e numéricas
            categoricas <- ['Categoria de geração']
            numericas <- ['Geração', 'Duração da Transição', 'Fase da Transição']
            
            // Configurar transformador de colunas
            transf <- ColumnTransformer([
                ('onehot', OneHotEncoder(sparse_output=FALSO, handle_unknown='ignore'), categoricas),
                ('scaler', StandardScaler(), numericas)
            ], remainder='passthrough', verbose_feature_names_out=FALSO)
            
            // Aplicar transformações
            data_tr <- transf.fit_transform(x_tr)
            data_val <- transf.transform(x_val)
            data_te <- transf.transform(x_te)
            
            // Criar DataFrames finais
            df_tr <- DataFrame(data_tr, colunas = transf.get_feature_names_out())
            df_val <- DataFrame(data_val, colunas = transf.get_feature_names_out())
            df_te <- DataFrame(data_te, colunas = transf.get_feature_names_out())
            
            // Extrair scaler para coluna 'Geração'
            scaler <- transf.named_transformers_['scaler']
            idx_pg <- posição de 'Geração' em numericas
            
            df_norm <- DataFrame({
                'Variável': ['Geração'],
                'Mean': [scaler.mean_[idx_pg]],
                'Std': [scaler.scale_[idx_pg]]
            })
            
            // Definir diretório de saída
            dir <- caminho('Usinas', usina, 'Minimização', pasta, 'Rede Neural', 'Dados')
            
            // Criar diretório se não existir
            criar_diretorio(dir, exist_ok=VERDADEIRO)
            
            // Salvar DataFrames
            salvar_parquet(df_tr, caminho(dir, 'Treino.parquet'), index=FALSO)
            salvar_parquet(df_val, caminho(dir, 'Validação.parquet'), index=FALSO)
            salvar_parquet(df_te, caminho(dir, 'Teste.parquet'), index=FALSO)
            salvar_parquet(df_norm, caminho(dir, 'Scaler Geração.parquet'), index=FALSO)
            
        FIM PARA
    FIM PARA

FIM ALGORITMO
\end{Verbatim}

\captionof{listing}{Treinamento da Rede Neural}
\begin{Verbatim}[breaklines=true, frame=single]
ALGORITMO treinar_modelo
ENTRADA: 
    model (rede neural)
    train_loader (dados de treino)
    val_loader (dados de validação)
    dyn_cols (colunas dinâmicas)
    epochs (inteiro, padrão = 100)
    patience (inteiro, padrão = 50)
    huber_delta (float, padrão = 0.5)
    peso_transicao (float, padrão = 8.0)
    clip_grad (float, padrão = 1.0)

SAÍDA:
    model (modelo treinado)
    history (dicionário com histórico de métricas)

INÍCIO
    // Mover modelo para dispositivo (CPU/GPU)
    model <- mover_para_device(model, device)
    
    // Configurar otimizador e funções de perda
    optimizer <- Adam(model.parameters(), lr=0.001)
    criterion_huber <- HuberLoss(delta=huber_delta)
    criterion_mae <- L1Loss()
    criterion_mse <- MSELoss()
    
    // Configurar scheduler para reduzir learning rate
    scheduler <- ReduceLROnPlateau(optimizer, mode='min', factor=0.5, 
                                   patience=5, min_lr=0.00001)
    
    // Inicializar variáveis de controle
    best_val_loss <- infinito
    patience_counter <- 0
    best_model_state <- None
    history <- {
        'loss': [], 'mae': [], 'mse': [],
        'val_loss': [], 'val_mae': [], 'val_mse': []
    }
    
    // Loop de épocas
    PARA epoch DE 1 ATÉ epochs:
        
        // === FASE DE TREINAMENTO ===
        model.treino()
        train_loss <- 0.0
        train_mae <- 0.0
        train_mse <- 0.0
        
        PARA CADA (x_dyn, x_sta, y) EM train_loader:
            // Mover dados para o dispositivo
            x_dyn <- mover_para_device(x_dyn, device)
            x_sta <- mover_para_device(x_sta, device)
            y <- mover_para_device(y, device)
            y <- squeeze(y)  // Remover dimensões extras
            
            // Zerar gradientes
            optimizer.zerar_gradientes()
            
            // Forward pass
            output <- model(x_dyn, x_sta)
            
            // Calcular perda base (Huber)
            loss_data <- criterion_huber(output, y)
            
            // Aplicar pesos para transições
            is_transicao <- (x_sta[:, 0] == 1)
            weight <- onde(is_transicao, peso_transicao, 1.0)
            loss <- media(loss_data * weight)
            
            // Backward pass
            loss.backward()
            
            // Clip gradients para evitar explosão
            clip_grad_norm_(model.parameters(), max_norm=clip_grad)
            
            // Atualizar pesos
            optimizer.step()
            
            // Acumular métricas
            train_loss <- train_loss + loss.item()
            train_mae <- train_mae + media(criterion_mae(output, y)).item()
            train_mse <- train_mse + media(criterion_mse(output, y)).item()
        
        // Calcular médias das métricas de treino
        train_loss <- train_loss / tamanho(train_loader)
        train_mae <- train_mae / tamanho(train_loader)
        train_mse <- train_mse / tamanho(train_loader)
        
        // === FASE DE VALIDAÇÃO ===
        model.avaliacao()  // Modo de avaliação (sem dropout, etc.)
        val_loss <- 0.0
        val_mae <- 0.0
        val_mse <- 0.0
        
        // Desabilitar cálculo de gradientes para validação
        COM torch.no_grad():
            PARA CADA (x_dyn, x_sta, y) EM val_loader:
                x_dyn <- mover_para_device(x_dyn, device)
                x_sta <- mover_para_device(x_sta, device)
                y <- mover_para_device(y, device)
                y <- squeeze(y)
                
                output <- model(x_dyn, x_sta)
                
                loss_val <- media(criterion_huber(output, y))
                val_loss <- val_loss + loss_val.item()
                val_mae <- val_mae + criterion_mae(output, y).item()
                val_mse <- val_mse + criterion_mse(output, y).item()
        
        // Calcular médias das métricas de validação
        val_loss <- val_loss / tamanho(val_loader)
        val_mae <- val_mae / tamanho(val_loader)
        val_mse <- val_mse / tamanho(val_loader)
        
        // Registrar histórico
        adicionar train_loss em history['loss']
        adicionar train_mae em history['mae']
        adicionar train_mse em history['mse']
        adicionar val_loss em history['val_loss']
        adicionar val_mae em history['val_mae']
        adicionar val_mse em history['val_mse']
        
        // Atualizar scheduler baseado na loss de validação
        scheduler.step(val_loss)
        
        // Verificar se é o melhor modelo até agora
        SE val_loss < best_val_loss ENTÃO
            best_val_loss <- val_loss
            patience_counter <- 0
            best_model_state <- copiar(model.state_dict())
        SENÃO
            patience_counter <- patience_counter + 1
            
            SE patience_counter >= patience ENTÃO
                INTERROMPER loop  // Early stopping
            FIM SE
        FIM SE
        
    FIM PARA
    
    // Carregar o melhor modelo encontrado
    SE best_model_state NÃO é None ENTÃO
        model.load_state_dict(best_model_state)
    FIM SE
    
    RETORNAR model, history

FIM ALGORITMO
\end{Verbatim}   % Arquivo apendice_1.tex
% \include{apendices/apendice_2}   % Arquivo apendice_2.tex
%%---------------------------------------------------------------

%%--------Anexos-------------------------------------------------
%% Descomente caso exista a necessidade da inclusão de anexos
%\anexo % Comando que inclui os anexos (outros autores)
\appendix
\chapter{Anexo}\label{chp:anexo}

  % Arquivo anexo_1.tex
% \include{anexos/anexo_2}  % Arquivo anexo_2.tex
%%---------------------------------------------------------------

\end{document}
